\documentclass[11pt]{article}
\usepackage{amsmath,amsthm,amssymb}
\usepackage[margin=1in]{geometry}
\usepackage[T1]{fontenc}

\newtheorem{theorem}{Theorem}
\newtheorem{lemma}{Lemma}
\newtheorem{corollary}{Corollary}
\newtheorem{proposition}{Proposition}

\begin{document}

\section*{Theorems}

\begin{theorem}[Theorem 5]
THEOREM 5. The space  $(S_m^+, \otimes)$  is an abelian Lie group. Moreover, the metric  $g$  defined in (3.1) is a bi-invariant metric on  $(S_m^+, \otimes)$ .
\end{theorem}

\begin{proof}
Proof. It is clear that  $\mathcal{S}_m^+$  is closed under the operation  $\circledast$ , and the identity element is the identity matrix. For  $P \in S_m^+$ , the inverse under  $\circledast$  is given by  $(\mathcal{L}(P)_{\odot}^{-1})(\mathcal{L}(P)_{\odot}^{-1})^\top$ . For associativity, we first observe that  $\mathcal{L}(P \circledast Q) = \mathcal{L}(P) \circledast \mathcal{L}(Q)$ , based on which we further deduce that

$$
\begin{array}{l} (P \circledast Q) \circledast S = (\mathcal {L} (P \circledast Q) \circledast \mathcal {L} (S)) (\mathcal {L} (P \circledast Q) \circledast \mathcal {L} (S)) ^ {\top} \\ = (\mathcal {L} (P) \circledast \mathcal {L} (Q) \circledast \mathcal {L} (S)) (\mathcal {L} (P) \circledast \mathcal {L} (Q) \circledast \mathcal {L} (S)) ^ {\top} \\ = (\mathcal {L} (P) \circledast \mathcal {L} (Q \circledast S)) (\mathcal {L} (P) \circledast \mathcal {L} (Q \circledast S)) ^ {\top} \\ = P \oplus (Q \oplus S). \\ \end{array}
$$

Therefore,  $(\mathcal{S}_m^+, \otimes)$  is a group. The commutativity and smoothness of  $\otimes$  stem from the commutativity and smoothness of  $\odot$ , respectively. It can be checked that  $\mathcal{S}$  is a group isomorphism and isometry between Lie groups  $(\mathcal{L}_+, \odot)$  and  $(\mathcal{S}_m^+, \otimes)$ , respectively, endowed with Riemannian metrics  $\tilde{g}$  and  $g$ . Then, the bi-invariance of  $g$  follows from the bi-invariance of  $\tilde{g}$ .

3.4. Parallel transport along geodesics on  $S_{m}^{+}$ . In some applications like statistical analysis or machine learning on Riemannian manifolds, parallel transport of tangent vectors along geodesics is required. For instance, in [40] that studies regression on SPD-valued data, tangent vectors are transported to a common place to derive statistical estimators of interest. Also, optimization in the context of statistics for manifold-valued data often involves parallel transport of tangent vectors. Examples include Hamiltonian Monte Carlo algorithms [23] as well as optimization algorithms [20] to train manifold-valued Gaussian mixture models. In these scenarios, Riemannian metrics on SPD matrices that result in efficient computation for parallel transport are attractive, in particular in the era of big data. In this regard, as discussed in the introduction, evaluation of parallel transport along geodesics for the affine-invariant metric is simple and efficient in computation, while the one for the Log-Euclidean metric is computationally intensive. Below we show that parallel transport for the presented metric also has a simple form, starting with a lemma.

\end{proof}

\paragraph{Background.}
$$
L ^ {- 1} W L ^ {- \top} = X ^ {\top} L ^ {- \top} + L ^ {- 1} X = L ^ {- 1} X + (L ^ {- 1} X) ^ {\top}.
$$

Note that  $L^{-1}X$  is also a lower triangular matrix, and the matrix on the left-hand side is symmetric, we deduce that  $(L^{-1}WL^{-\top})_{\frac{1}{2}} = L^{-1}X$ , which gives rise to  $X = L(L^{-1}WL^{-\top})_{\frac{1}{2}}$ . The linear map  $(D_L\mathcal{S})(X) = LX^\top +XL^\top$  is one-to-one, since from the above derivation,  $(D_L\mathcal{S})(X) = 0$  if and only if  $X = 0$ .

Given the above proposition, the manifold map  $\mathcal{S}$ , which is exactly the inverse of the Cholesky map  $\mathcal{L}$  discussed in subsection 2.2, induces a Riemannian metric on  $S_{m}^{+}$ , denoted by  $g$  and called Log-Cholesky metric, given by

Note that  $L^{-1}X$  is also a lower triangular matrix, and the matrix on the left-hand side is symmetric, we deduce that  $(L^{-1}WL^{-\top})_{\frac{1}{2}} = L^{-1}X$ , which gives rise to  $X = L(L^{-1}WL^{-\top})_{\frac{1}{2}}$ . The linear map  $(D_L\mathcal{S})(X) = LX^\top +XL^\top$  is one-to-one, since from the above derivation,  $(D_L\mathcal{S})(X) = 0$  if and only if  $X = 0$ .

Given the above proposition, the manifold map  $\mathcal{S}$ , which is exactly the inverse of the Cholesky map  $\mathcal{L}$  discussed in subsection 2.2, induces a Riemannian metric on  $S_{m}^{+}$ , denoted by  $g$  and called Log-Cholesky metric, given by

$$
g _ {P} (W, V) = \tilde {g} _ {L} \left(L \left(L ^ {- 1} W L ^ {- \top}\right) _ {\frac {1}{2}}, L \left(L ^ {- 1} V L ^ {- \top}\right) _ {\frac {1}{2}}\right), \tag {3.1}
$$

The group operator  $\odot$  and maps  $\mathcal{S}$  and  $\mathcal{L}$  together induce a smooth operation  $\otimes$  on  $S_{m}^{+}$ , defined by

$$
\begin{array}{l} P \circledast Q = \mathcal {S} (\mathcal {L} (P) \circledast \mathcal {L} (Q)) \\ = \left(\mathcal {L} (P) \circledast \mathcal {L} (Q)\right) \left(\mathcal {L} (P) \circledast \mathcal {L} (Q)\right) ^ {\top} \quad \text {f o r} P, Q \in S _ {m} ^ {+}. \\ \end{array}
$$

In addition, both  $\mathcal{L}$  and  $\mathcal{S}$  are Riemannian isometries and group isomorphisms between Lie groups  $(\mathcal{L}_{+},\tilde{g},\circledast)$  and  $(\mathcal{S}_m^+,g,\bullet)$ .

$$
\begin{array}{l} P \circledast Q = \mathcal {S} (\mathcal {L} (P) \circledast \mathcal {L} (Q)) \\ = \left(\mathcal {L} (P) \circledast \mathcal {L} (Q)\right) \left(\mathcal {L} (P) \circledast \mathcal {L} (Q)\right) ^ {\top} \quad \text {f o r} P, Q \in S _ {m} ^ {+}. \\ \end{array}
$$

In addition, both  $\mathcal{L}$  and  $\mathcal{S}$  are Riemannian isometries and group isomorphisms between Lie groups  $(\mathcal{L}_{+},\tilde{g},\circledast)$  and  $(\mathcal{S}_m^+,g,\bullet)$ .

THEOREM 5. The space  $(S_m^+, \otimes)$  is an abelian Lie group. Moreover, the metric  $g$  defined in (3.1) is a bi-invariant metric on  $(S_m^+, \otimes)$ .


\section*{Lemmas}

\begin{lemma}[Lemma 6]
LEMMA 6. Let  $(\mathcal{G},\cdot)$  be an abelian Lie group with a bi-invariant metric. The parallel transporter  $\tau_{p,q}$  that transports tangent vectors at  $p$  to tangent vectors at  $q$  along the geodesic connecting  $p$  and  $q$  is given by  $\tau_{p,q}(u) = (D_p\ell_{q,p^{-1}})u$  for  $u\in T_p\mathcal{G}$ .
\end{lemma}

\begin{proof}
Proof. For simplicity, we abbreviate  $p \cdot q$  as  $pq$ . Let  $\mathfrak{g}$  denote the Lie algebra associated with the Lie group  $\mathcal{G}$ , and  $\nabla$  the Levi-Civita connection on  $\mathcal{G}$ . Note that we identify elements in  $\mathfrak{g}$  with left-invariant vector fields on  $\mathcal{G}$ . We shall first recall that  $\nabla_{Y}Z = [Y,Z] / 2$  for  $Y,Z \in \mathfrak{g}$  (see the proof of Theorem 21.3 in [29]), where  $[\cdot ,\cdot ]$  denotes the Lie bracket of  $\mathcal{G}$ . As  $\mathcal{G}$  is abelian, the Lie bracket vanishes everywhere, and hence  $\nabla_{Y}Z = 0$  if  $Y,Z \in \mathfrak{g}$ .

Let  $\gamma_p(t) = \ell_p(\mathfrak{exp}(tY))$  for  $Y \in \mathfrak{g}$  such that  $\mathfrak{exp}(Y) = p^{-1}q$ , where  $\mathfrak{exp}$  denotes the Lie group exponential map. Recall that for a bi-invariant Lie group, the group exponential map coincides with the Riemannian exponential map at the group identity  $e$ , and left translations are isometries. Thus,  $\gamma_p$  is a geodesic. Using the fact  $\gamma_e(t + s) = \gamma_e(t)\gamma_e(s) = \ell_{\gamma_e(t)}(\gamma_e(s))$  according to Lemma 21.2 of [29], by the chain rule of differential, we have

$$
\gamma_ {e} ^ {\prime} (t) = \frac {\mathrm {d}}{\mathrm {d} t} \gamma_ {e} (t + s) \mid_ {s = 0} = (D _ {e} \ell_ {\gamma_ {e} (t)}) (\gamma_ {e} ^ {\prime} (0)) = (D _ {e} \ell_ {\gamma_ {e} (t)}) (Y (e)) = Y (\gamma_ {e} (t)),
$$

from which we further deduce that

$$
\gamma_ {p} ^ {\prime} (t) = (D _ {\gamma_ {e} (t)} \ell_ {p}) \gamma_ {e} ^ {\prime} (t) = (D _ {\gamma_ {e} (t)} \ell_ {p}) Y (\gamma_ {e} (t)) = Y (p \gamma_ {e} (t)).
$$

Now define a vector field  $Z(q) \coloneqq (D_p\ell_{qp^{-1}})u$ . We claim that  $Z$  is a left-invariant vector field on  $\mathcal{G}$  and hence belongs to  $\mathfrak{g}$ , since

$$
\begin{array}{l} Z (h q) = \left(D _ {p} \ell_ {h q p ^ {- 1}}\right) u = \left\{D _ {p} \left(\ell_ {h} \circ \ell_ {q p ^ {- 1}}\right) \right\} u \\ = \left(D _ {q} \ell_ {h}\right) \left(D _ {p} \ell_ {q p ^ {- 1}}\right) u = \left(D _ {q} \ell_ {h}\right) (Z (q)), \\ \end{array}
$$

where the third equality is obtained by the chain rule of differential. Consequently,  $\nabla_{\gamma_p'}Z = 0$  since  $\gamma_p'(t) = Y(p\gamma_e(t))$  and  $\nabla_Y Z = 0$  for  $Y,Z\in \mathfrak{g}$ . As additionally  $Z(\gamma_p(0)) = Z(p) = u$ , transportation of  $u$  along the geodesic  $\gamma_{p}$  is realized by the left-invariant vector field  $Z$ . Since  $\gamma_{p}$  is a geodesic with  $\gamma_p(0) = p$  and  $\gamma_p(1) = \ell_p\mathfrak{e}\mathfrak{p}(Y) = p(p^{-1}q) = q$ , we have that

$$
\tau_ {p, q} (u) = Z \left(\gamma_ {p} (1)\right) = Z (q) = \left(D _ {p} \ell_ {q p ^ {- 1}}\right) u,
$$

as claimed.

\end{proof}

\paragraph{Background.}
$$
\begin{array}{l} Z (h q) = \left(D _ {p} \ell_ {h q p ^ {- 1}}\right) u = \left\{D _ {p} \left(\ell_ {h} \circ \ell_ {q p ^ {- 1}}\right) \right\} u \\ = \left(D _ {q} \ell_ {h}\right) \left(D _ {p} \ell_ {q p ^ {- 1}}\right) u = \left(D _ {q} \ell_ {h}\right) (Z (q)), \\ \end{array}
$$

where the third equality is obtained by the chain rule of differential. Consequently,  $\nabla_{\gamma_p'}Z = 0$  since  $\gamma_p'(t) = Y(p\gamma_e(t))$  and  $\nabla_Y Z = 0$  for  $Y,Z\in \mathfrak{g}$ . As additionally  $Z(\gamma_p(0)) = Z(p) = u$ , transportation of  $u$  along the geodesic  $\gamma_{p}$  is realized by the left-invariant vector field  $Z$ . Since  $\gamma_{p}$  is a geodesic with  $\gamma_p(0) = p$  and  $\gamma_p(1) = \ell_p\mathfrak{e}\mathfrak{p}(Y) = p(p^{-1}q) = q$ , we have that

$$
\tau_ {p, q} (u) = Z \left(\gamma_ {p} (1)\right) = Z (q) = \left(D _ {p} \ell_ {q p ^ {- 1}}\right) u,
$$

where the third equality is obtained by the chain rule of differential. Consequently,  $\nabla_{\gamma_p'}Z = 0$  since  $\gamma_p'(t) = Y(p\gamma_e(t))$  and  $\nabla_Y Z = 0$  for  $Y,Z\in \mathfrak{g}$ . As additionally  $Z(\gamma_p(0)) = Z(p) = u$ , transportation of  $u$  along the geodesic  $\gamma_{p}$  is realized by the left-invariant vector field  $Z$ . Since  $\gamma_{p}$  is a geodesic with  $\gamma_p(0) = p$  and  $\gamma_p(1) = \ell_p\mathfrak{e}\mathfrak{p}(Y) = p(p^{-1}q) = q$ , we have that

$$
\tau_ {p, q} (u) = Z \left(\gamma_ {p} (1)\right) = Z (q) = \left(D _ {p} \ell_ {q p ^ {- 1}}\right) u,
$$

as claimed.


\section*{Propositions}

\begin{proposition}[Proposition 2]
PROPOSITION 2. The Cholesky map  $\mathcal{L}$  is a diffeomorphism between smooth manifolds  $\mathcal{L}_{+}$  and  $S_{m}^{+}$ .
\end{proposition}

\begin{proof}
Proof. We have argued that  $\mathcal{L}$  is a bijection. To see that it is also smooth, for  $P = LL^{\top}$  with  $L\in \mathcal{L}_{+}$ , we write

$$
\begin{array}{l} \left( \begin{array}{c c c c} P _ {1 1} & P _ {1 2} & \dots & P _ {1 m} \\ P _ {2 1} & P _ {2 2} & \dots & P _ {2 m} \\ \vdots & \vdots & \ddots & \vdots \\ P _ {m 1} & P _ {m 2} & \dots & P _ {m m} \end{array} \right) = \left( \begin{array}{c c c c} L _ {1 1} & 0 & \dots & 0 \\ L _ {2 1} & L _ {2 2} & \dots & 0 \\ \vdots & \vdots & \ddots & \vdots \\ L _ {m 1} & L _ {m 2} & \dots & L _ {m m} \end{array} \right) \left( \begin{array}{c c c c} L _ {1 1} & L _ {2 1} & \dots & L _ {m 1} \\ 0 & L _ {2 2} & \dots & L _ {m 2} \\ \vdots & \vdots & \ddots & \vdots \\ 0 & 0 & \dots & L _ {m m} \end{array} \right) \\ = \left( \begin{array}{c c c c} L _ {1 1} ^ {2} & L _ {2 1} L _ {1 1} & \dots & L _ {m 1} L _ {1 1} \\ L _ {2 1} L _ {1 1} & L _ {2 1} ^ {2} + L _ {2 2} ^ {2} & \dots & L _ {m 1} L _ {2 1} + L _ {m 2} L _ {2 2} \\ \vdots & \vdots & \ddots & \vdots \\ L _ {m 1} L _ {1 1} & L _ {m 1} L _ {2 1} + L _ {m 2} L _ {2 2} & \dots & \sum_ {k = 1} ^ {m} L _ {m k} ^ {2} \end{array} \right), \\ \end{array}
$$

from which we deduce that

$$
\left\{ \begin{array}{l} L _ {i i} = \sqrt {P _ {i i} - \sum_ {k = 1} ^ {i - 1} L _ {i k} ^ {2}}, \\ L _ {i j} = \frac {1}{L _ {j j}} \left(P _ {i j} - \sum_ {k = 1} ^ {j - 1} L _ {i k} L _ {j k}\right) \quad \text {f o r} i > j. \end{array} \right. \tag {2.1}
$$

The existence of a unique Cholesky factor for every SPD matrix suggests that  $P_{ii} - \sum_{k=1}^{i-1} L_{ik}^2 > 0$  for all  $i$ . Thus,  $L_{11} = \sqrt{P_{11}}$ , as well as its reciprocal  $1 / L_{11}$ , is smooth.

Now assume  $L_{ij}$  and  $1 / L_{jj}$  are smooth for  $i = 1, \ldots, i_0$  and  $j = 1, \ldots, j_0 \leq i_0$ . As we just showed, this hypothesis is true for  $i_0 = 1$  and  $j_0 = 1$ . If  $j_0 = i_0$ , from (2.1) we see that  $L_{i_0 + 1,1} = (1 / L_{11})P_{i_0 + 1,1}$  is smooth. If  $j_0 < i_0 - 1$ , then  $L_{i_0,j_0 + 1}$  results from a sequence of elementary operations, such as multiplication, addition, and subtraction, of maps  $L_{i_01}, \ldots, L_{i_0,j_0}, L_{j_0 + 1,1}, \ldots, L_{j_0 + 1,j_0}$ , and  $1 / L_{j_0 + 1,j_0 + 1}$  that are all smooth according to the induction hypothesis. As these elementary operations are all smooth,  $L_{i_0,j_0 + 1}$  is also smooth. If  $j_0 = i_0 - 1$ , then  $L_{i_0,j_0 + 1} = L_{i_0,i_0}$ , as well as  $1 / L_{i_0,i_0}$ , is smooth via similar reasoning based on the additional fact that  $P_{i_0,i_0} - \sum_{k=1}^{i_0-1}L_{ik}^2 > 0$  and the square-root operator  $\sqrt{}$  is smooth on the set of positive real numbers. The above derivation then shows that the induction hypothesis is also true for  $i = i_0, j = j_0 + 1$  if  $j_0 < i_0$  and  $i = i_0 + 1, j = 1$  if  $j_0 = i_0$ .

Consequently, by mathematical induction, the hypothesis is true for all pairs of  $i$  and  $j \leq i$ . In other words,  $\mathcal{L}$  is a smooth manifold map. Its inverse, denoted by  $\mathcal{S}$ , is given by  $\mathcal{S}(L) = LL^{\top}$  and clearly smooth. Therefore,  $\mathcal{L}$  and its inverse  $\mathcal{S}$  are diffeomorphisms.

3. Lie group structure and bi-invariant metric. In this section, we first construct a group structure and a bi-invariant metric on the manifold  $\mathcal{L}_{+}$  and then push them forward to the manifold  $S_{m}^{+}$  via the Cholesky map. Parallel transport on SPD manifolds is also investigated. For a background of Riemannian geometry and Lie group, we recommend monographs [19, 24].

3.1. Riemannian geometry on Cholesky spaces. For matrices in  $\mathcal{L}_{+}$ , as off-diagonal elements in the lower triangular part are unconstrained while diagonal ones are restricted to be positive,  $\mathcal{L}_{+}$  can be parameterized by  $\mathcal{L} \ni X \coloneqq \varphi(X) \in \mathcal{L}_{+}$  in the way that  $(\varphi(X))_{ij} = X_{ij}$  if  $i \neq j$  and  $(\varphi(X))_{jj} = \exp(X_{jj})$ . This motivates us to, respectively, endow the unconstrained part  $\lfloor \cdot \rfloor$  and the positive diagonal part  $\mathbb{D}(\cdot)$  with a different metric and then combine them into a Riemannian metric on  $\mathcal{L}_{+}$  as follows. First, we note that the tangent space of  $\mathcal{L}_{+}$  at a given  $L \in \mathcal{L}_{+}$  is identified with the linear space  $\mathcal{L}$ . For such tangent space, we treat the strict lower triangular space  $\lfloor \mathcal{L} \rfloor \coloneqq \{\lfloor X \rfloor : X \in \mathcal{L}\}$  as the Euclidean space  $\mathbb{R}^{m(m-1)/2}$  with the usual Frobenius inner product  $\langle X, Y \rangle_{\mathrm{F}} = \sum_{i,j=1}^{m} X_{ij} Y_{ij}$  for all  $X, Y \in \lfloor \mathcal{L} \rfloor$ . For the diagonal part  $\mathbb{D}(\mathcal{L}) \coloneqq \{\mathbb{D}(Z) : Z \in \mathcal{L}\}$ , we equipped it with a different inner product defined by  $\langle \mathbb{D}(L)^{-1} \mathbb{D}(X), \mathbb{D}(L)^{-1} \mathbb{D}(Y) \rangle_{\mathrm{F}}$ . Finally, combining these two components together, we define a metric  $\tilde{g}$  for tangent spaces  $T_L \mathcal{L}_{+}$  (identified with  $\mathcal{L}$ ) by

$$
\begin{array}{l} \tilde {g} _ {L} (X, Y) = \langle \lfloor X \rfloor , \lfloor Y \rfloor \rangle_ {\mathrm {F}} + \langle \mathbb {D} (L) ^ {- 1} \mathbb {D} (X), \mathbb {D} (L) ^ {- 1} \mathbb {D} (Y) \rangle_ {\mathrm {F}} \\ = \sum_ {i > j} X _ {i j} Y _ {i j} + \sum_ {j = 1} ^ {m} X _ {j j} Y _ {j j} L _ {j j} ^ {- 2}. \\ \end{array}
$$

It is straightforward to show that the space  $\mathcal{L}_{+}$ , equipped with the metric  $\tilde{g}$ , is a Riemannian manifold. We begin with geodesics on the manifold  $(\mathcal{L}_{+},\tilde{g})$  to investigate its basic properties.

\end{proof}

\paragraph{Background.}
-  $\operatorname{det}(X) = \prod_{j=1}^{m} X_{jj}$  for  $X \in \mathcal{L}$ .

These properties show that both  $\mathcal{L}$  and  $\mathcal{L}_{+}$  are closed under matrix addition and multiplication and that the operator  $\mathbb{D}$  interacts well with these operations.

Recall that  $S_{m}^{+}$  is defined as the collection of  $m \times m$  SPD matrices. We denote the space of  $m \times m$  symmetric space by  $S_{m}$ . Symmetric matrices and SPD matrices possess numerous algebraic and analytic properties that are well documented in [6]. Below are some of them to be used in what follows.

-  $\exp(S)$  is an SPD matrix for a symmetric matrix  $S$ , while  $\log(P)$  is a symmetric metric for an SPD matrix  $P$ .

- Diagonal elements of an SPD matrix are all positive. This can be seen from the fact that  $P_{jj} = e_j^\top Pe_j > 0$  for  $P \in S_m^+$ , where  $e_j$  is the unit vector with 1 at the  $j$ th coordinate and 0 elsewhere.

2.2. Cholesky decomposition. Cholesky decomposition, named after André-Louis Cholesky, represents a real  $m \times m$  SPD matrix  $P$  as a product of a lower triangular matrix  $L$  and its transpose, i.e.,  $P = LL^{\top}$ . If the diagonal elements of  $L$  are restricted to be positive, then the decomposition is unique according to Theorem 4.2.5 of [16]. Such lower triangular matrix, denoted by  $\mathcal{L}(P)$ , is called the Cholesky factor of  $P$ . Since in addition  $L = \mathcal{L}(LL^{\top})$  for each  $L \in \mathcal{L}_{+}$ , the map  $\mathcal{L}: S_{m}^{+} \to \mathcal{L}_{+}$  is bijective. In other words, there is one-to-one correspondence between SPD matrices and lower triangular matrices whose diagonal elements are all positive.

3.1. Riemannian geometry on Cholesky spaces. For matrices in  $\mathcal{L}_{+}$ , as off-diagonal elements in the lower triangular part are unconstrained while diagonal ones are restricted to be positive,  $\mathcal{L}_{+}$  can be parameterized by  $\mathcal{L} \ni X \coloneqq \varphi(X) \in \mathcal{L}_{+}$  in the way that  $(\varphi(X))_{ij} = X_{ij}$  if  $i \neq j$  and  $(\varphi(X))_{jj} = \exp(X_{jj})$ . This motivates us to, respectively, endow the unconstrained part  $\lfloor \cdot \rfloor$  and the positive diagonal part  $\mathbb{D}(\cdot)$  with a different metric and then combine them into a Riemannian metric on  $\mathcal{L}_{+}$  as follows. First, we note that the tangent space of  $\mathcal{L}_{+}$  at a given  $L \in \mathcal{L}_{+}$  is identified with the linear space  $\mathcal{L}$ . For such tangent space, we treat the strict lower triangular space  $\lfloor \mathcal{L} \rfloor \coloneqq \{\lfloor X \rfloor : X \in \mathcal{L}\}$  as the Euclidean space  $\mathbb{R}^{m(m-1)/2}$  with the usual Frobenius inner product  $\langle X, Y \rangle_{\mathrm{F}} = \sum_{i,j=1}^{m} X_{ij} Y_{ij}$  for all  $X, Y \in \lfloor \mathcal{L} \rfloor$ . For the diagonal part  $\mathbb{D}(\mathcal{L}) \coloneqq \{\mathbb{D}(Z) : Z \in \mathcal{L}\}$ , we equipped it with a different inner product defined by  $\langle \mathbb{D}(L)^{-1} \mathbb{D}(X), \mathbb{D}(L)^{-1} \mathbb{D}(Y) \rangle_{\mathrm{F}}$ . Finally, combining these two components together, we define a metric  $\tilde{g}$  for tangent spaces  $T_L \mathcal{L}_{+}$  (identified with  $\mathcal{L}$ ) by

$$
\begin{array}{l} \tilde {g} _ {L} (X, Y) = \langle \lfloor X \rfloor , \lfloor Y \rfloor \rangle_ {\mathrm {F}} + \langle \mathbb {D} (L) ^ {- 1} \mathbb {D} (X), \mathbb {D} (L) ^ {- 1} \mathbb {D} (Y) \rangle_ {\mathrm {F}} \\ = \sum_ {i > j} X _ {i j} Y _ {i j} + \sum_ {j = 1} ^ {m} X _ {j j} Y _ {j j} L _ {j j} ^ {- 2}. \\ \end{array}
$$

It is straightforward to show that the space  $\mathcal{L}_{+}$ , equipped with the metric  $\tilde{g}$ , is a Riemannian manifold. We begin with geodesics on the manifold  $(\mathcal{L}_{+},\tilde{g})$  to investigate its basic properties.

$$
L ^ {- 1} W L ^ {- \top} = X ^ {\top} L ^ {- \top} + L ^ {- 1} X = L ^ {- 1} X + (L ^ {- 1} X) ^ {\top}.
$$

Note that  $L^{-1}X$  is also a lower triangular matrix, and the matrix on the left-hand side is symmetric, we deduce that  $(L^{-1}WL^{-\top})_{\frac{1}{2}} = L^{-1}X$ , which gives rise to  $X = L(L^{-1}WL^{-\top})_{\frac{1}{2}}$ . The linear map  $(D_L\mathcal{S})(X) = LX^\top +XL^\top$  is one-to-one, since from the above derivation,  $(D_L\mathcal{S})(X) = 0$  if and only if  $X = 0$ .

Given the above proposition, the manifold map  $\mathcal{S}$ , which is exactly the inverse of the Cholesky map  $\mathcal{L}$  discussed in subsection 2.2, induces a Riemannian metric on  $S_{m}^{+}$ , denoted by  $g$  and called Log-Cholesky metric, given by

Note that  $L^{-1}X$  is also a lower triangular matrix, and the matrix on the left-hand side is symmetric, we deduce that  $(L^{-1}WL^{-\top})_{\frac{1}{2}} = L^{-1}X$ , which gives rise to  $X = L(L^{-1}WL^{-\top})_{\frac{1}{2}}$ . The linear map  $(D_L\mathcal{S})(X) = LX^\top +XL^\top$  is one-to-one, since from the above derivation,  $(D_L\mathcal{S})(X) = 0$  if and only if  $X = 0$ .

Given the above proposition, the manifold map  $\mathcal{S}$ , which is exactly the inverse of the Cholesky map  $\mathcal{L}$  discussed in subsection 2.2, induces a Riemannian metric on  $S_{m}^{+}$ , denoted by  $g$  and called Log-Cholesky metric, given by

$$
g _ {P} (W, V) = \tilde {g} _ {L} \left(L \left(L ^ {- 1} W L ^ {- \top}\right) _ {\frac {1}{2}}, L \left(L ^ {- 1} V L ^ {- \top}\right) _ {\frac {1}{2}}\right), \tag {3.1}
$$

3.3. Lie group structure and bi-invariant metrics. We define an operator  $\odot$  on  $\mathcal{L}$  by

$$
X \circledast Y = \left\lfloor X \right\rfloor + \left\lfloor Y \right\rfloor + \mathbb {D} (X) \mathbb {D} (Y).
$$

Note that  $\mathcal{L}_{+} \subset \mathcal{L}$ . Moreover, if  $L, K \in \mathcal{L}_{+}$ , then  $L \circledast K \in \mathcal{L}_{+}$ . It is not difficult to see that  $\circledast$  is a smooth commutative group operation on the manifold  $\mathcal{L}_{+}$ , where the

$$
X \circledast Y = \left\lfloor X \right\rfloor + \left\lfloor Y \right\rfloor + \mathbb {D} (X) \mathbb {D} (Y).
$$

Note that  $\mathcal{L}_{+} \subset \mathcal{L}$ . Moreover, if  $L, K \in \mathcal{L}_{+}$ , then  $L \circledast K \in \mathcal{L}_{+}$ . It is not difficult to see that  $\circledast$  is a smooth commutative group operation on the manifold  $\mathcal{L}_{+}$ , where the

# LOG-CHOLESKY METRIC

inverse of  $L$ , denoted by  $L_{\odot}^{-1}$ , is  $\mathbb{D}(L)^{-1} - \lfloor L\rfloor$ . The left translation by  $A\in \mathcal{L}_{+}$  is denoted by  $\ell_A:B\mapsto A\odot B$ . One can check that the differential of this operation at  $L\in \mathcal{L}_{+}$  is

$$
D _ {L} \ell_ {A}: X \mapsto \lfloor X \rfloor + \mathbb {D} (A) \mathbb {D} (X), \tag {3.2}
$$

where it is noted that the differential  $D_{L}\ell_{A}$  does not depend on  $L$ . Given the above expression, one can find that

$$
D _ {L} \ell_ {A}: X \mapsto \lfloor X \rfloor + \mathbb {D} (A) \mathbb {D} (X), \tag {3.2}
$$

where it is noted that the differential  $D_{L}\ell_{A}$  does not depend on  $L$ . Given the above expression, one can find that

$$
\begin{array}{l} \tilde {g} _ {A \odot L} ((D _ {L} \ell_ {A}) (X), (D _ {L} \ell_ {A}) (Y)) \\ = \tilde {g} _ {A \odot L} (\left\lfloor X \right\rfloor + \mathbb {D} (A) \mathbb {D} (X), \left\lfloor Y \right\rfloor + \mathbb {D} (A) \mathbb {D} (Y)) \\ = \tilde {g} _ {L} (X, Y). \\ \end{array}
$$

The group operator  $\odot$  and maps  $\mathcal{S}$  and  $\mathcal{L}$  together induce a smooth operation  $\otimes$  on  $S_{m}^{+}$ , defined by

$$
\begin{array}{l} P \circledast Q = \mathcal {S} (\mathcal {L} (P) \circledast \mathcal {L} (Q)) \\ = \left(\mathcal {L} (P) \circledast \mathcal {L} (Q)\right) \left(\mathcal {L} (P) \circledast \mathcal {L} (Q)\right) ^ {\top} \quad \text {f o r} P, Q \in S _ {m} ^ {+}. \\ \end{array}
$$

In addition, both  $\mathcal{L}$  and  $\mathcal{S}$  are Riemannian isometries and group isomorphisms between Lie groups  $(\mathcal{L}_{+},\tilde{g},\circledast)$  and  $(\mathcal{S}_m^+,g,\bullet)$ .

$$
\begin{array}{l} P \circledast Q = \mathcal {S} (\mathcal {L} (P) \circledast \mathcal {L} (Q)) \\ = \left(\mathcal {L} (P) \circledast \mathcal {L} (Q)\right) \left(\mathcal {L} (P) \circledast \mathcal {L} (Q)\right) ^ {\top} \quad \text {f o r} P, Q \in S _ {m} ^ {+}. \\ \end{array}
$$

In addition, both  $\mathcal{L}$  and  $\mathcal{S}$  are Riemannian isometries and group isomorphisms between Lie groups  $(\mathcal{L}_{+},\tilde{g},\circledast)$  and  $(\mathcal{S}_m^+,g,\bullet)$ .

THEOREM 5. The space  $(S_m^+, \otimes)$  is an abelian Lie group. Moreover, the metric  $g$  defined in (3.1) is a bi-invariant metric on  $(S_m^+, \otimes)$ .


\begin{proposition}[Proposition 3]
PROPOSITION 3. On the Riemannian manifold  $(\mathcal{L}_{+},\tilde{g})$ , the geodesic starting at  $L\in \mathcal{L}_{+}$  with direction  $X\in T_{L}\mathcal{L}_{+}$  is given by

$$
\tilde {\gamma} _ {L, X} (t) = \left\lfloor L \right\rfloor + t \left\lfloor X \right\rfloor + \mathbb {D} (L) \exp \{t \mathbb {D} (X) \mathbb {D} (L) ^ {- 1} \}.
$$
\end{proposition}

\begin{proof}
Proof. Clearly,  $\tilde{\gamma}_{L,X}(0) = L$  and  $\tilde{\gamma}_{L,X}'(0) = \lfloor X\rfloor +\mathbb{D}(X) = X$ . Now, we use  $\operatorname{vec}(L)$  to denote the vector in  $\mathbb{R}^{m(m+1)/2}$  such that the first  $m$  elements of  $\operatorname{vec}(L)$  correspond to the diagonal elements of  $L$ . Define the map  $x:\mathcal{L}_+\to \mathbb{R}$  by

$$
x ^ {i} (L) = \left\{ \begin{array}{l l} \log \operatorname {v e c} (L) _ {i} & \text {i f} 1 \leq i \leq m, \\ \operatorname {v e c} (L) _ {i}, & \text {o t h e r w i s e}, \end{array} \right.
$$

where  $x^{i}$  denotes the  $i$ th component of  $x$ . It can be checked that  $(\mathcal{L}_{+},x)$  is a chart for the manifold  $\mathcal{L}_{+}$ . Let  $e_{i}$  be the  $m(m + 1) / 2$  dimensional vector whose  $i$ th element is one and other elements are all zero. For  $1\leq i\leq m$ , we define  $\partial_i = \operatorname {vec}(L)_ie_i$ , and for  $i > m$ , define  $\partial_i = e_i$ . The collection  $\{\partial_1,\dots ,\partial_{m(m + 1) / 2}\}$  is a frame. One can check that  $\tilde{g}_{ij}\coloneqq \tilde{g}_L(\partial_i,\partial_j) = 0$  if  $i\neq j$ , and  $\tilde{g}_{ii} = 1$ . This implies that  $\partial \tilde{g}_{jk} / \partial x^l = 0$ , and hence all Christoffel symbols are zero, as

$$
\Gamma_ {k l} ^ {i} = \frac {1}{2} \tilde {g} ^ {i j} \left(\frac {\partial \tilde {g} _ {j k}}{\partial x ^ {l}} + \frac {\partial \tilde {g} _ {j l}}{\partial x ^ {k}} - \frac {\partial \tilde {g} _ {k l}}{\partial x ^ {j}}\right) = 0,
$$

where Einstein summation convention is assumed. It can be checked that the  $i$ th coordinate  $\tilde{\gamma}^i(t) = x^i \circ \tilde{\gamma}_{L,X}(t)$  of the curve  $\tilde{\gamma}_{L,X}$  is given by  $\tilde{\gamma}^i(t) = \log \operatorname{vec}(L)_i + t\operatorname{vec}(X)_i / \operatorname{vec}(L)_i$  when  $i \leq m$  and  $\tilde{\gamma}^i(t) = \operatorname{vec}(L)_i + t\operatorname{vec}(X)_i$  if  $i > m$ . Now, it is an easy task to verify the following geodesic equations

$$
\frac {d ^ {2} \tilde {\gamma} ^ {i}}{d t ^ {2}} + \Gamma_ {j k} ^ {i} \frac {d \tilde {\gamma} ^ {j}}{d t} \frac {d \tilde {\gamma} ^ {k}}{d t} = 0
$$

for  $i = 1, \dots, m(m + 1) / 2$ . Therefore,  $\tilde{\gamma}_{L,X}(t)$  is the claimed geodesic.

Given the above proposition, we can immediately derive the Riemannian exponential map  $\operatorname{Exp}$  at  $L \in \mathcal{L}_{+}$ , which is given by

$$
\widetilde {\mathrm {E x p}} _ {L} X = \widetilde {\gamma} _ {L, X} (1) = \left\lfloor L \right\rfloor + \left\lfloor X \right\rfloor + \mathbb {D} (L) \exp \{\mathbb {D} (X) \mathbb {D} (L) ^ {- 1} \}.
$$

Also, for  $L, K \in \mathcal{L}_{+}$ , with  $X = \lfloor K \rfloor - \lfloor L \rfloor + \{\log \mathbb{D}(K) - \log \mathbb{D}(L)\} \mathbb{D}(L)$ , one has

$$
\tilde {\gamma} _ {L, X} (t) = \lfloor L \rfloor + t \{\lfloor K \rfloor - \lfloor L \rfloor \} + \mathbb {D} (L) \exp \left[ t \left\{\log \mathbb {D} (K) - \log \mathbb {D} (L) \right\} \right].
$$

Since  $\tilde{\gamma}_{L,X}(0) = L$  and  $\tilde{\gamma}_{L,X}(1) = K$ ,  $\tilde{\gamma}_{L,X}$  is the geodesic connecting  $L$  and  $K$ . Therefore, the distance function on  $\mathcal{L}_+$  induced by  $\tilde{g}$ , denoted by  $d_{\mathcal{L}_+}$ , is given by

$$
d _ {\mathcal {L} _ {+}} (L, K) = \{\tilde {g} _ {L} (X, X) \} ^ {1 / 2} = \left\{\sum_ {i > j} (L _ {i j} - K _ {i j}) ^ {2} + \sum_ {j = 1} ^ {m} (\log L _ {j j} - \log K _ {j j}) ^ {2} \right\} ^ {1 / 2},
$$

where  $X$  is the same as the above. The expression for the distance function can be equivalently and more compactly written as

$$
d _ {\mathcal {L} _ {+}} (L, K) = \{\| \lfloor L \rfloor - \lfloor K \rfloor \| _ {\mathrm {F}} ^ {2} + \| \log \mathbb {D} (L) - \log \mathbb {D} (K) \| _ {\mathrm {F}} ^ {2} \} ^ {1 / 2}.
$$

Table 1 summarizes the above basic properties of the manifold  $(\mathcal{L}_{+},\tilde{g})$


TABLE 1 Basic properties of Riemannian manifolds  $(\mathcal{L}_{+},\tilde{g})$


<table><tr><td>tangent space at L
L</td></tr><tr><td>Riemannian metric
gL(X,Y) = ∑i&gt;jXijYij + ∑m
j=1XjjYjjL-2j</td></tr><tr><td>geodesic emanating from L with direction X
γL,X(t) = [L] + t[X] + D(L) exp{tD(X)D(L)-1}</td></tr><tr><td>Riemannian exponential map at L
ExpL X = [L] + [X] + D(L) exp{D(X)D(L)-1}</td></tr><tr><td>Riemannian logarithmic map at L
LogL K = [K] - [L] + D(L) log{D(L)-1D(K)}</td></tr><tr><td>geodesic distance between L and K
dL_+(L,K) = {||[L] - [K]||_F^2 + ||logD(L) - logD(K)||_F^2}^{1/2}</td></tr></table>

3.2. Riemannian metric for SPD matrices. As mentioned previously, the space  $S_{m}^{+}$  of SPD matrices is a smooth submanifold of the space  $S_{m}$  of symmetric matrices, whose tangent space at a given SPD matrix is identified with  $S_{m}$ . We also showed that the map  $\mathcal{S}:\mathcal{L}_{+}\to S_{m}^{+}$  by  $\mathcal{S}(L) = LL^{\top}$  is a diffeomorphism between  $\mathcal{L}_{+}$  and  $S_{m}^{+}$ . For a square matrix  $S$ , we define a lower triangular matrix  $S_{\frac{1}{2}} = \lfloor S\rfloor +\mathbb{D}(S) / 2$ . In other words, the matrix  $S_{\frac{1}{2}}$  is the lower triangular part of  $S$  with the diagonal elements halved. The differential of  $\mathcal{S}$  is given in the following.

\end{proof}

\paragraph{Background.}
- For any invertible matrix  $X$ , the matrix  $XX^{\top}$  is an SPD matrix.

-  $\exp(S)$  is an SPD matrix for a symmetric matrix  $S$ , while  $\log(P)$  is a symmetric metric for an SPD matrix  $P$ .

- Diagonal elements of an SPD matrix are all positive. This can be seen from the fact that  $P_{jj} = e_j^\top Pe_j > 0$  for  $P \in S_m^+$ , where  $e_j$  is the unit vector with 1 at the  $j$ th coordinate and 0 elsewhere.

-  $\exp(S)$  is an SPD matrix for a symmetric matrix  $S$ , while  $\log(P)$  is a symmetric metric for an SPD matrix  $P$ .

- Diagonal elements of an SPD matrix are all positive. This can be seen from the fact that  $P_{jj} = e_j^\top Pe_j > 0$  for  $P \in S_m^+$ , where  $e_j$  is the unit vector with 1 at the  $j$ th coordinate and 0 elsewhere.

2.2. Cholesky decomposition. Cholesky decomposition, named after André-Louis Cholesky, represents a real  $m \times m$  SPD matrix  $P$  as a product of a lower triangular matrix  $L$  and its transpose, i.e.,  $P = LL^{\top}$ . If the diagonal elements of  $L$  are restricted to be positive, then the decomposition is unique according to Theorem 4.2.5 of [16]. Such lower triangular matrix, denoted by  $\mathcal{L}(P)$ , is called the Cholesky factor of  $P$ . Since in addition  $L = \mathcal{L}(LL^{\top})$  for each  $L \in \mathcal{L}_{+}$ , the map  $\mathcal{L}: S_{m}^{+} \to \mathcal{L}_{+}$  is bijective. In other words, there is one-to-one correspondence between SPD matrices and lower triangular matrices whose diagonal elements are all positive.

3.1. Riemannian geometry on Cholesky spaces. For matrices in  $\mathcal{L}_{+}$ , as off-diagonal elements in the lower triangular part are unconstrained while diagonal ones are restricted to be positive,  $\mathcal{L}_{+}$  can be parameterized by  $\mathcal{L} \ni X \coloneqq \varphi(X) \in \mathcal{L}_{+}$  in the way that  $(\varphi(X))_{ij} = X_{ij}$  if  $i \neq j$  and  $(\varphi(X))_{jj} = \exp(X_{jj})$ . This motivates us to, respectively, endow the unconstrained part  $\lfloor \cdot \rfloor$  and the positive diagonal part  $\mathbb{D}(\cdot)$  with a different metric and then combine them into a Riemannian metric on  $\mathcal{L}_{+}$  as follows. First, we note that the tangent space of  $\mathcal{L}_{+}$  at a given  $L \in \mathcal{L}_{+}$  is identified with the linear space  $\mathcal{L}$ . For such tangent space, we treat the strict lower triangular space  $\lfloor \mathcal{L} \rfloor \coloneqq \{\lfloor X \rfloor : X \in \mathcal{L}\}$  as the Euclidean space  $\mathbb{R}^{m(m-1)/2}$  with the usual Frobenius inner product  $\langle X, Y \rangle_{\mathrm{F}} = \sum_{i,j=1}^{m} X_{ij} Y_{ij}$  for all  $X, Y \in \lfloor \mathcal{L} \rfloor$ . For the diagonal part  $\mathbb{D}(\mathcal{L}) \coloneqq \{\mathbb{D}(Z) : Z \in \mathcal{L}\}$ , we equipped it with a different inner product defined by  $\langle \mathbb{D}(L)^{-1} \mathbb{D}(X), \mathbb{D}(L)^{-1} \mathbb{D}(Y) \rangle_{\mathrm{F}}$ . Finally, combining these two components together, we define a metric  $\tilde{g}$  for tangent spaces  $T_L \mathcal{L}_{+}$  (identified with  $\mathcal{L}$ ) by

$$
\begin{array}{l} \tilde {g} _ {L} (X, Y) = \langle \lfloor X \rfloor , \lfloor Y \rfloor \rangle_ {\mathrm {F}} + \langle \mathbb {D} (L) ^ {- 1} \mathbb {D} (X), \mathbb {D} (L) ^ {- 1} \mathbb {D} (Y) \rangle_ {\mathrm {F}} \\ = \sum_ {i > j} X _ {i j} Y _ {i j} + \sum_ {j = 1} ^ {m} X _ {j j} Y _ {j j} L _ {j j} ^ {- 2}. \\ \end{array}
$$

It is straightforward to show that the space  $\mathcal{L}_{+}$ , equipped with the metric  $\tilde{g}$ , is a Riemannian manifold. We begin with geodesics on the manifold  $(\mathcal{L}_{+},\tilde{g})$  to investigate its basic properties.

$$
d _ {\mathcal {L} _ {+}} (L, K) = \{\tilde {g} _ {L} (X, X) \} ^ {1 / 2} = \left\{\sum_ {i > j} (L _ {i j} - K _ {i j}) ^ {2} + \sum_ {j = 1} ^ {m} (\log L _ {j j} - \log K _ {j j}) ^ {2} \right\} ^ {1 / 2},
$$

where  $X$  is the same as the above. The expression for the distance function can be equivalently and more compactly written as

$$
d _ {\mathcal {L} _ {+}} (L, K) = \{\| \lfloor L \rfloor - \lfloor K \rfloor \| _ {\mathrm {F}} ^ {2} + \| \log \mathbb {D} (L) - \log \mathbb {D} (K) \| _ {\mathrm {F}} ^ {2} \} ^ {1 / 2}.
$$

$$
L ^ {- 1} W L ^ {- \top} = X ^ {\top} L ^ {- \top} + L ^ {- 1} X = L ^ {- 1} X + (L ^ {- 1} X) ^ {\top}.
$$

Note that  $L^{-1}X$  is also a lower triangular matrix, and the matrix on the left-hand side is symmetric, we deduce that  $(L^{-1}WL^{-\top})_{\frac{1}{2}} = L^{-1}X$ , which gives rise to  $X = L(L^{-1}WL^{-\top})_{\frac{1}{2}}$ . The linear map  $(D_L\mathcal{S})(X) = LX^\top +XL^\top$  is one-to-one, since from the above derivation,  $(D_L\mathcal{S})(X) = 0$  if and only if  $X = 0$ .

Given the above proposition, the manifold map  $\mathcal{S}$ , which is exactly the inverse of the Cholesky map  $\mathcal{L}$  discussed in subsection 2.2, induces a Riemannian metric on  $S_{m}^{+}$ , denoted by  $g$  and called Log-Cholesky metric, given by

Note that  $L^{-1}X$  is also a lower triangular matrix, and the matrix on the left-hand side is symmetric, we deduce that  $(L^{-1}WL^{-\top})_{\frac{1}{2}} = L^{-1}X$ , which gives rise to  $X = L(L^{-1}WL^{-\top})_{\frac{1}{2}}$ . The linear map  $(D_L\mathcal{S})(X) = LX^\top +XL^\top$  is one-to-one, since from the above derivation,  $(D_L\mathcal{S})(X) = 0$  if and only if  $X = 0$ .

Given the above proposition, the manifold map  $\mathcal{S}$ , which is exactly the inverse of the Cholesky map  $\mathcal{L}$  discussed in subsection 2.2, induces a Riemannian metric on  $S_{m}^{+}$ , denoted by  $g$  and called Log-Cholesky metric, given by

$$
g _ {P} (W, V) = \tilde {g} _ {L} \left(L \left(L ^ {- 1} W L ^ {- \top}\right) _ {\frac {1}{2}}, L \left(L ^ {- 1} V L ^ {- \top}\right) _ {\frac {1}{2}}\right), \tag {3.1}
$$

inverse of  $L$ , denoted by  $L_{\odot}^{-1}$ , is  $\mathbb{D}(L)^{-1} - \lfloor L\rfloor$ . The left translation by  $A\in \mathcal{L}_{+}$  is denoted by  $\ell_A:B\mapsto A\odot B$ . One can check that the differential of this operation at  $L\in \mathcal{L}_{+}$  is

$$
D _ {L} \ell_ {A}: X \mapsto \lfloor X \rfloor + \mathbb {D} (A) \mathbb {D} (X), \tag {3.2}
$$

where it is noted that the differential  $D_{L}\ell_{A}$  does not depend on  $L$ . Given the above expression, one can find that

$$
D _ {L} \ell_ {A}: X \mapsto \lfloor X \rfloor + \mathbb {D} (A) \mathbb {D} (X), \tag {3.2}
$$

where it is noted that the differential  $D_{L}\ell_{A}$  does not depend on  $L$ . Given the above expression, one can find that

$$
\begin{array}{l} \tilde {g} _ {A \odot L} ((D _ {L} \ell_ {A}) (X), (D _ {L} \ell_ {A}) (Y)) \\ = \tilde {g} _ {A \odot L} (\left\lfloor X \right\rfloor + \mathbb {D} (A) \mathbb {D} (X), \left\lfloor Y \right\rfloor + \mathbb {D} (A) \mathbb {D} (Y)) \\ = \tilde {g} _ {L} (X, Y). \\ \end{array}
$$

$$
\begin{array}{l} P \circledast Q = \mathcal {S} (\mathcal {L} (P) \circledast \mathcal {L} (Q)) \\ = \left(\mathcal {L} (P) \circledast \mathcal {L} (Q)\right) \left(\mathcal {L} (P) \circledast \mathcal {L} (Q)\right) ^ {\top} \quad \text {f o r} P, Q \in S _ {m} ^ {+}. \\ \end{array}
$$

In addition, both  $\mathcal{L}$  and  $\mathcal{S}$  are Riemannian isometries and group isomorphisms between Lie groups  $(\mathcal{L}_{+},\tilde{g},\circledast)$  and  $(\mathcal{S}_m^+,g,\bullet)$ .

THEOREM 5. The space  $(S_m^+, \otimes)$  is an abelian Lie group. Moreover, the metric  $g$  defined in (3.1) is a bi-invariant metric on  $(S_m^+, \otimes)$ .

$$
\begin{array}{l} \det  \mathbb {E} L = \exp [ \operatorname {t r} \{\log \mathbb {D} (\mathbb {E} L) \} ] = \exp [ \operatorname {t r} \{\mathbb {E} \log \mathbb {D} (L) \} ] \\ = \exp \left[ \mathbb {E} \operatorname {t r} \left\{\log \mathbb {D} (L) \right\} \right] = \exp \left\{\mathbb {E} \log (\det  L) \right\}, \\ \end{array}
$$

which proves the case of  $(\mathcal{L}_{+},\tilde{g})$

For the case that  $Q$  is a random element on  $(\mathcal{S}_m^+, g)$ , we denote it by  $P$  and observe that  $\operatorname{det}(AB) = (\operatorname{det} A)(\operatorname{det} B)$  for any square matrices  $A$  and  $B$ . Let  $L = \mathcal{L}(P)$ . Then, one has

which proves the case of  $(\mathcal{L}_{+},\tilde{g})$

For the case that  $Q$  is a random element on  $(\mathcal{S}_m^+, g)$ , we denote it by  $P$  and observe that  $\operatorname{det}(AB) = (\operatorname{det} A)(\operatorname{det} B)$  for any square matrices  $A$  and  $B$ . Let  $L = \mathcal{L}(P)$ . Then, one has

$$
\begin{array}{l} \det  \mathbb {E} P = \det  \left\{\left(\mathbb {E} L\right) \left(\mathbb {E} L\right) ^ {\top} \right\} \\ = \left\{\det  (\mathbb {E} L) \right\} \left\{\det  (\mathbb {E} L) ^ {\top} \right\} \\ = \exp \left\{\mathbb {E} \log (\det  L) \right\} \exp \left\{\mathbb {E} \log (\det  L ^ {\top}) \right\} \\ = \exp \left\{\mathbb {E} \left(\log \det  L + \log \det  L ^ {\top}\right) \right\} \\ = \exp \left[ \mathbb {E} \log \left\{\left(\det  L\right) \det  \left(L ^ {\top}\right) \right\} \right], \\ \end{array}
$$


\begin{proposition}[Proposition 4]
PROPOSITION 4. The differential  $D_L\mathcal{S}:T_L\mathcal{L}_+ \to T_{LL^\top}\mathcal{S}_m^+$  of  $\mathcal{S}$  at  $L$  is given by

$$
\left(D _ {L} \mathcal {S}\right) (X) = L X ^ {\top} + X L ^ {\top}.
$$
\end{proposition}

\begin{proof}
Proof. Let  $X \in \mathcal{L}$  and  $L \in \mathcal{L}_{+}$ . Then  $\tilde{\gamma}_L(t) = L + tX$  is a curve passing through  $L$  if  $t \in (-\epsilon, \epsilon)$  for a sufficiently small  $\epsilon > 0$ . Note that for every such  $t$ ,  $\tilde{\gamma}_L(t) \in \mathcal{L}_{+}$ . Then  $\gamma_{LL^{\top}}(t) = \mathcal{S}(\tilde{\gamma}_L(t))$  is a curve passing through  $LL^{\top}$ . The differential is then derived from

$$
\left(D _ {L} \mathcal {S}\right) (X) = \left(\mathcal {S} \circ \tilde {\gamma} _ {L}\right) ^ {\prime} (0) = \frac {\mathrm {d}}{\mathrm {d} t} \mathcal {S} \left(\tilde {\gamma} _ {L} (t)\right) | _ {t = 0} = L X ^ {\top} + X L ^ {\top}.
$$

On the other hand, if  $W = LX^{\top} + XL^{\top}$ , then since  $L$  is invertible, we have

$$
L ^ {- 1} W L ^ {- \top} = X ^ {\top} L ^ {- \top} + L ^ {- 1} X = L ^ {- 1} X + (L ^ {- 1} X) ^ {\top}.
$$

Note that  $L^{-1}X$  is also a lower triangular matrix, and the matrix on the left-hand side is symmetric, we deduce that  $(L^{-1}WL^{-\top})_{\frac{1}{2}} = L^{-1}X$ , which gives rise to  $X = L(L^{-1}WL^{-\top})_{\frac{1}{2}}$ . The linear map  $(D_L\mathcal{S})(X) = LX^\top +XL^\top$  is one-to-one, since from the above derivation,  $(D_L\mathcal{S})(X) = 0$  if and only if  $X = 0$ .

Given the above proposition, the manifold map  $\mathcal{S}$ , which is exactly the inverse of the Cholesky map  $\mathcal{L}$  discussed in subsection 2.2, induces a Riemannian metric on  $S_{m}^{+}$ , denoted by  $g$  and called Log-Cholesky metric, given by

$$
g _ {P} (W, V) = \tilde {g} _ {L} \left(L \left(L ^ {- 1} W L ^ {- \top}\right) _ {\frac {1}{2}}, L \left(L ^ {- 1} V L ^ {- \top}\right) _ {\frac {1}{2}}\right), \tag {3.1}
$$

where  $L = \mathcal{S}^{-1}(P) = \mathcal{L}(P)$  is the Cholesky factor of  $P \in S_m^+$ , and  $W, V \in S_m$  are tangent vectors at  $P$ . This implies that

$$
\tilde {g} _ {L} (X, Y) = g _ {\mathcal {S} (L)} \big ((D _ {L} \mathcal {S}) (X), (D _ {L} \mathcal {S}) (Y) \big)
$$

for all  $L$  and  $X,Y\in T_L\mathcal{L}_+$ . According to Definition 7.57 of [26], the map  $\mathcal{S}$  is an isometry between  $(\mathcal{L}_{+},\tilde{g})$  and  $(S_m^+,g)$ . A Riemannian isometry provides correspondence of Riemannian properties and objects between two Riemannian manifolds. This enables us to study the properties of  $(S_m^+,g)$  via the manifold  $(\mathcal{L}_{+},\tilde{g})$  and the isometry  $\mathcal{S}$ . For example, we can obtain geodesics on  $S_{m}^{+}$  by mapping geodesics on  $\mathcal{L}_{+}$ . More precisely, the geodesic emanating from  $P = LL^{\top}$  with  $L = \mathcal{L}(P)$  is given by

$$
\gamma_ {P, W} (t) = \mathcal {S} (\tilde {\gamma} _ {L, X} (t)) = \tilde {\gamma} _ {L, X} (t) \tilde {\gamma} _ {L, X} ^ {\top} (t),
$$

where  $X = L(L^{-1}WL^{-\top})_{\frac{1}{2}}$  and  $W \in T_P\mathcal{S}_m^+$ . Similarly, the Riemannian exponential at  $P$  is given by

$$
\mathrm {E x p} _ {P} W = \mathcal {S} (\widetilde {\mathrm {E x p}} _ {L} X) = (\widetilde {\mathrm {E x p}} _ {L} X) (\widetilde {\mathrm {E x p}} _ {L} X) ^ {\top},
$$

while the geodesic between  $P$  and  $Q$  is characterized by

$$
\gamma_ {P, W} (t) = \tilde {\gamma} _ {L, X} (t) \tilde {\gamma} _ {L, X} ^ {\top} (t),
$$

with  $L = \mathcal{L}(P)$ ,  $K = \mathcal{L}(Q)$ ,  $X = \lfloor K\rfloor -\lfloor L\rfloor +\{\log \mathbb{D}(K) - \log \mathbb{D}(L)\} \mathbb{D}(L)$ , and  $W = LX^{\top} + XL^{\top}$ . Also, the geodesic distance between  $P$  and  $Q$  is

$$
d _ {\mathcal {S} _ {m} ^ {+}} (P, Q) = d _ {\mathcal {L} _ {+}} (\mathcal {L} (P), \mathcal {L} (Q)).
$$

Moreover, the Levi-Civita connection  $\nabla$  of  $(\mathcal{S}_m^+,g)$  can be obtained by the Levi-Civita connection  $\tilde{\nabla}$  of  $(\mathcal{L}_{+},\tilde{g})$ . To see this, let  $W$  and  $V$  be two smooth vector fields on  $\mathcal{S}_m^+$ . Define vector fields  $X$  and  $Y$  on  $\mathcal{L}_{+}$  by  $X(L) = (D_{LL^{\top}}\mathcal{L})W(LL^{\top})$  and  $Y(L) = (D_{LL^{\top}}\mathcal{L})V(LL^{\top})$  for all  $L\in \mathcal{L}_{+}$ . Then  $\nabla_W V = (D\mathcal{S})(\tilde{\nabla}_XY)$ , and the Christoffel symbols to compute the connection  $\tilde{\nabla}$  has been given in the proof of Proposition 3.

Table 2 summarizes some basic properties of the manifold  $(S_m^+, g)$ . Note that the differential  $D_P\mathcal{L}$  can be computed efficiently, since it only involves Cholesky decomposition and the inverse of a lower triangular matrix, for both of which there exist efficient algorithms. Consequently, all maps in Table 2 can be evaluated in an efficient way. In contrast, computation of Riemannian exponential/logarithmic maps for the Log-Euclidean metric [4] requires evaluation of some series of an infinite number of terms; see (2.1) and Table 4.1 of [4].

3.3. Lie group structure and bi-invariant metrics. We define an operator  $\odot$  on  $\mathcal{L}$  by

$$
X \circledast Y = \left\lfloor X \right\rfloor + \left\lfloor Y \right\rfloor + \mathbb {D} (X) \mathbb {D} (Y).
$$

Note that  $\mathcal{L}_{+} \subset \mathcal{L}$ . Moreover, if  $L, K \in \mathcal{L}_{+}$ , then  $L \circledast K \in \mathcal{L}_{+}$ . It is not difficult to see that  $\circledast$  is a smooth commutative group operation on the manifold  $\mathcal{L}_{+}$ , where the

\end{proof}

\paragraph{Background.}
- For any invertible matrix  $X$ , the matrix  $XX^{\top}$  is an SPD matrix.

-  $\exp(S)$  is an SPD matrix for a symmetric matrix  $S$ , while  $\log(P)$  is a symmetric metric for an SPD matrix  $P$ .

- Diagonal elements of an SPD matrix are all positive. This can be seen from the fact that  $P_{jj} = e_j^\top Pe_j > 0$  for  $P \in S_m^+$ , where  $e_j$  is the unit vector with 1 at the  $j$ th coordinate and 0 elsewhere.

-  $\exp(S)$  is an SPD matrix for a symmetric matrix  $S$ , while  $\log(P)$  is a symmetric metric for an SPD matrix  $P$ .

- Diagonal elements of an SPD matrix are all positive. This can be seen from the fact that  $P_{jj} = e_j^\top Pe_j > 0$  for  $P \in S_m^+$ , where  $e_j$  is the unit vector with 1 at the  $j$ th coordinate and 0 elsewhere.

2.2. Cholesky decomposition. Cholesky decomposition, named after André-Louis Cholesky, represents a real  $m \times m$  SPD matrix  $P$  as a product of a lower triangular matrix  $L$  and its transpose, i.e.,  $P = LL^{\top}$ . If the diagonal elements of  $L$  are restricted to be positive, then the decomposition is unique according to Theorem 4.2.5 of [16]. Such lower triangular matrix, denoted by  $\mathcal{L}(P)$ , is called the Cholesky factor of  $P$ . Since in addition  $L = \mathcal{L}(LL^{\top})$  for each  $L \in \mathcal{L}_{+}$ , the map  $\mathcal{L}: S_{m}^{+} \to \mathcal{L}_{+}$  is bijective. In other words, there is one-to-one correspondence between SPD matrices and lower triangular matrices whose diagonal elements are all positive.

$$
d _ {\mathcal {L} _ {+}} (L, K) = \{\tilde {g} _ {L} (X, X) \} ^ {1 / 2} = \left\{\sum_ {i > j} (L _ {i j} - K _ {i j}) ^ {2} + \sum_ {j = 1} ^ {m} (\log L _ {j j} - \log K _ {j j}) ^ {2} \right\} ^ {1 / 2},
$$

where  $X$  is the same as the above. The expression for the distance function can be equivalently and more compactly written as

$$
d _ {\mathcal {L} _ {+}} (L, K) = \{\| \lfloor L \rfloor - \lfloor K \rfloor \| _ {\mathrm {F}} ^ {2} + \| \log \mathbb {D} (L) - \log \mathbb {D} (K) \| _ {\mathrm {F}} ^ {2} \} ^ {1 / 2}.
$$

$$
L ^ {- 1} W L ^ {- \top} = X ^ {\top} L ^ {- \top} + L ^ {- 1} X = L ^ {- 1} X + (L ^ {- 1} X) ^ {\top}.
$$

Note that  $L^{-1}X$  is also a lower triangular matrix, and the matrix on the left-hand side is symmetric, we deduce that  $(L^{-1}WL^{-\top})_{\frac{1}{2}} = L^{-1}X$ , which gives rise to  $X = L(L^{-1}WL^{-\top})_{\frac{1}{2}}$ . The linear map  $(D_L\mathcal{S})(X) = LX^\top +XL^\top$  is one-to-one, since from the above derivation,  $(D_L\mathcal{S})(X) = 0$  if and only if  $X = 0$ .

Given the above proposition, the manifold map  $\mathcal{S}$ , which is exactly the inverse of the Cholesky map  $\mathcal{L}$  discussed in subsection 2.2, induces a Riemannian metric on  $S_{m}^{+}$ , denoted by  $g$  and called Log-Cholesky metric, given by

Note that  $L^{-1}X$  is also a lower triangular matrix, and the matrix on the left-hand side is symmetric, we deduce that  $(L^{-1}WL^{-\top})_{\frac{1}{2}} = L^{-1}X$ , which gives rise to  $X = L(L^{-1}WL^{-\top})_{\frac{1}{2}}$ . The linear map  $(D_L\mathcal{S})(X) = LX^\top +XL^\top$  is one-to-one, since from the above derivation,  $(D_L\mathcal{S})(X) = 0$  if and only if  $X = 0$ .

Given the above proposition, the manifold map  $\mathcal{S}$ , which is exactly the inverse of the Cholesky map  $\mathcal{L}$  discussed in subsection 2.2, induces a Riemannian metric on  $S_{m}^{+}$ , denoted by  $g$  and called Log-Cholesky metric, given by

$$
g _ {P} (W, V) = \tilde {g} _ {L} \left(L \left(L ^ {- 1} W L ^ {- \top}\right) _ {\frac {1}{2}}, L \left(L ^ {- 1} V L ^ {- \top}\right) _ {\frac {1}{2}}\right), \tag {3.1}
$$

inverse of  $L$ , denoted by  $L_{\odot}^{-1}$ , is  $\mathbb{D}(L)^{-1} - \lfloor L\rfloor$ . The left translation by  $A\in \mathcal{L}_{+}$  is denoted by  $\ell_A:B\mapsto A\odot B$ . One can check that the differential of this operation at  $L\in \mathcal{L}_{+}$  is

$$
D _ {L} \ell_ {A}: X \mapsto \lfloor X \rfloor + \mathbb {D} (A) \mathbb {D} (X), \tag {3.2}
$$

where it is noted that the differential  $D_{L}\ell_{A}$  does not depend on  $L$ . Given the above expression, one can find that

$$
D _ {L} \ell_ {A}: X \mapsto \lfloor X \rfloor + \mathbb {D} (A) \mathbb {D} (X), \tag {3.2}
$$

where it is noted that the differential  $D_{L}\ell_{A}$  does not depend on  $L$ . Given the above expression, one can find that

$$
\begin{array}{l} \tilde {g} _ {A \odot L} ((D _ {L} \ell_ {A}) (X), (D _ {L} \ell_ {A}) (Y)) \\ = \tilde {g} _ {A \odot L} (\left\lfloor X \right\rfloor + \mathbb {D} (A) \mathbb {D} (X), \left\lfloor Y \right\rfloor + \mathbb {D} (A) \mathbb {D} (Y)) \\ = \tilde {g} _ {L} (X, Y). \\ \end{array}
$$

The group operator  $\odot$  and maps  $\mathcal{S}$  and  $\mathcal{L}$  together induce a smooth operation  $\otimes$  on  $S_{m}^{+}$ , defined by

$$
\begin{array}{l} P \circledast Q = \mathcal {S} (\mathcal {L} (P) \circledast \mathcal {L} (Q)) \\ = \left(\mathcal {L} (P) \circledast \mathcal {L} (Q)\right) \left(\mathcal {L} (P) \circledast \mathcal {L} (Q)\right) ^ {\top} \quad \text {f o r} P, Q \in S _ {m} ^ {+}. \\ \end{array}
$$

In addition, both  $\mathcal{L}$  and  $\mathcal{S}$  are Riemannian isometries and group isomorphisms between Lie groups  $(\mathcal{L}_{+},\tilde{g},\circledast)$  and  $(\mathcal{S}_m^+,g,\bullet)$ .

$$
\begin{array}{l} P \circledast Q = \mathcal {S} (\mathcal {L} (P) \circledast \mathcal {L} (Q)) \\ = \left(\mathcal {L} (P) \circledast \mathcal {L} (Q)\right) \left(\mathcal {L} (P) \circledast \mathcal {L} (Q)\right) ^ {\top} \quad \text {f o r} P, Q \in S _ {m} ^ {+}. \\ \end{array}
$$

In addition, both  $\mathcal{L}$  and  $\mathcal{S}$  are Riemannian isometries and group isomorphisms between Lie groups  $(\mathcal{L}_{+},\tilde{g},\circledast)$  and  $(\mathcal{S}_m^+,g,\bullet)$ .

THEOREM 5. The space  $(S_m^+, \otimes)$  is an abelian Lie group. Moreover, the metric  $g$  defined in (3.1) is a bi-invariant metric on  $(S_m^+, \otimes)$ .

[37] S. SRA, Positive definite matrices and the  $S$ -divergence, Proc. Amer. Math. Soc., 144 (2016), pp. 2787-2797.

[38] B. VANDEREYCKEN, P.-A. ABSIL, AND S. VANDEWALLE, A Riemannian geometry with complete geodesics for the set of positive semidefinite matrices of fixed rank, IMA J. Numer. Anal., 33 (2013), pp. 481-514.

[39] Z. WANG, B. C. VEMURI, Y. CHEN, AND T. H. MARECI, A constrained variational principle for direct estimation and smoothing of the diffusion tensor field from complex DWI, IEEE Trans. Medical Imaging, 23 (2004), pp. 930-939.


\begin{proposition}[Proposition 7]
PROPOSITION 7. A tangent vector  $X \in T_{L}\mathcal{L}_{+}$  is parallelly transported to the tangent vector  $\lfloor X\rfloor + \mathbb{D}(K)\mathbb{D}(L)^{-1}\mathbb{D}(X)$  at  $K$ . For  $P, Q \in S_m^+$  and  $W \in S_m$ ,

$$
\tau_ {P, Q} (W) = K \{\lfloor X \rfloor + \mathbb {D} (K) \mathbb {D} (L) ^ {- 1} \mathbb {D} (X) \} ^ {\top} + \{\lfloor X \rfloor + \mathbb {D} (K) \mathbb {D} (L ^ {- 1}) \mathbb {D} (X) \} K ^ {\top},
$$

where  $L = \mathcal{L}(P)$ ,  $K = \mathcal{L}(Q)$  and  $X = L(L^{-1}WL^{-\top})_{\frac{1}{2}}$ .
\end{proposition}

\begin{proof}
Proof. By Lemma 6, it is seen that  $\tau_{L,K}(X) = (D_L\ell_{K\mathbb{O}L_{\mathbb{O}}^{-1}})X$ . According to (3.2),

$$
(D _ {L} \ell_ {K \circledast L _ {\circledast} ^ {- 1}}) X = \lfloor X \rfloor + \mathbb {D} (K \circledast L _ {\circledast} ^ {- 1}) \mathbb {D} (X) = \lfloor X \rfloor + \mathbb {D} (K) \mathbb {D} (L) ^ {- 1} \mathbb {D} (X).
$$

The statement for  $P, Q, W$  follows from the fact that  $\mathcal{S}$  and  $\mathcal{L}$  are an isometries.  $\square$

The above proposition shows that parallel transport for the presented Log-Cholesky metric can be computed rather efficiently. In fact, the only nontrivial steps in computation are to perform Cholesky decomposition of two SPD matrices and to invert a lower triangular matrix, for both of which there exist efficient algorithms. It is numerically faster than the affine-invariant metric and Log-Euclidean metric. For instance, for  $5 \times 5$  SPD matrices, in a MATLAB computational environment, on average it takes  $9.3\mathrm{ms}$  (Log-Euclidean),  $0.85\mathrm{ms}$  (affine-invariant) and  $0.2\mathrm{ms}$  (Log-Cholesky) on an Intel(R) Core i7-4500U (1.80 GHz) to perform a parallel transport. We see that, to do parallel transport, the Log-Euclidean metric is about 45 times slower than the Log-Cholesky metric, since it has to evaluate the differential of a matrix logarithmic map that is expressed as a convergent series of infinite terms of products of matrices. In practice, only a finite leading terms are evaluated. However, in order to avoid significant loss of precision, a large number of terms are often needed. For instance, for a precision of the order of  $10^{-12}$ , typically approximately 300 leading terms are required. The affine-invariant metric, despite having a simple form for parallel transport, is still about 4 times slower than our Log-Cholesky metric, partially due to that more matrix inversion and multiplication operations are needed.

4. Mean of distributions on  $S_{m}^{+}$ . In this section we study the Log-Cholesky mean of a random SPD matrix and the Log-Cholesky average of a finite number of SPD matrices. We first establish the existence and uniqueness of such quantities. A closed and easy-to-compute form for Log-Cholesky averages is also derived. Finally, we show that determinants of Log-Cholesky averages are bounded by determinants of SPD matrices being averaged. This property suggests that Log-Cholesky averages are not subject to swelling effect.

4.1. Log-Cholesky mean of a random SPD matrix. For a random element  $Q$  on a Riemannian manifold  $\mathcal{M}$ , we define a function  $F(x) = \mathbb{E}d_{\mathcal{M}}^2 (x,Q)$ , where  $d_{\mathcal{M}}$  denotes the geodesic distance function on  $\mathcal{M}$ , and  $\mathbb{E}$  denotes expectation of a random number. If  $F(x) < \infty$  for some  $x\in \mathcal{M}$ ,

$$
x = \operatorname * {a r g   m i n} _ {z \in \mathcal {M}} F (z),
$$

and  $F(z) \geq F(x)$  for all  $z \in \mathcal{M}$ , then  $x$  is called a Fréchet mean of  $Q$ , denoted by  $\mathbb{E}Q$ . In general, Fréchet mean might not exist, and when it exists, it might not be unique; see [1] for conditions of the existence and uniqueness of Fréchet means. However, for  $S_{m}^{+}$  endowed with Log-Cholesky metric, we claim that if  $F(S) < \infty$  for some  $S \in S_{m}^{+}$ , then the Fréchet mean exists and is unique. Such Fréchet mean is termed Log-Cholesky mean in this paper. To prove the claim, we first notice that, similar to the Log-Euclidean metric [4], the manifold  $(S_{m}^{+}, g)$  is a "flat" space.

\end{proof}

\paragraph{Background.}
-  $\operatorname{det}(X) = \prod_{j=1}^{m} X_{jj}$  for  $X \in \mathcal{L}$ .

These properties show that both  $\mathcal{L}$  and  $\mathcal{L}_{+}$  are closed under matrix addition and multiplication and that the operator  $\mathbb{D}$  interacts well with these operations.

Recall that  $S_{m}^{+}$  is defined as the collection of  $m \times m$  SPD matrices. We denote the space of  $m \times m$  symmetric space by  $S_{m}$ . Symmetric matrices and SPD matrices possess numerous algebraic and analytic properties that are well documented in [6]. Below are some of them to be used in what follows.

- For any invertible matrix  $X$ , the matrix  $XX^{\top}$  is an SPD matrix.

-  $\exp(S)$  is an SPD matrix for a symmetric matrix  $S$ , while  $\log(P)$  is a symmetric metric for an SPD matrix  $P$ .

- Diagonal elements of an SPD matrix are all positive. This can be seen from the fact that  $P_{jj} = e_j^\top Pe_j > 0$  for  $P \in S_m^+$ , where  $e_j$  is the unit vector with 1 at the  $j$ th coordinate and 0 elsewhere.

-  $\exp(S)$  is an SPD matrix for a symmetric matrix  $S$ , while  $\log(P)$  is a symmetric metric for an SPD matrix  $P$ .

- Diagonal elements of an SPD matrix are all positive. This can be seen from the fact that  $P_{jj} = e_j^\top Pe_j > 0$  for  $P \in S_m^+$ , where  $e_j$  is the unit vector with 1 at the  $j$ th coordinate and 0 elsewhere.

2.2. Cholesky decomposition. Cholesky decomposition, named after André-Louis Cholesky, represents a real  $m \times m$  SPD matrix  $P$  as a product of a lower triangular matrix  $L$  and its transpose, i.e.,  $P = LL^{\top}$ . If the diagonal elements of  $L$  are restricted to be positive, then the decomposition is unique according to Theorem 4.2.5 of [16]. Such lower triangular matrix, denoted by  $\mathcal{L}(P)$ , is called the Cholesky factor of  $P$ . Since in addition  $L = \mathcal{L}(LL^{\top})$  for each  $L \in \mathcal{L}_{+}$ , the map  $\mathcal{L}: S_{m}^{+} \to \mathcal{L}_{+}$  is bijective. In other words, there is one-to-one correspondence between SPD matrices and lower triangular matrices whose diagonal elements are all positive.

3.1. Riemannian geometry on Cholesky spaces. For matrices in  $\mathcal{L}_{+}$ , as off-diagonal elements in the lower triangular part are unconstrained while diagonal ones are restricted to be positive,  $\mathcal{L}_{+}$  can be parameterized by  $\mathcal{L} \ni X \coloneqq \varphi(X) \in \mathcal{L}_{+}$  in the way that  $(\varphi(X))_{ij} = X_{ij}$  if  $i \neq j$  and  $(\varphi(X))_{jj} = \exp(X_{jj})$ . This motivates us to, respectively, endow the unconstrained part  $\lfloor \cdot \rfloor$  and the positive diagonal part  $\mathbb{D}(\cdot)$  with a different metric and then combine them into a Riemannian metric on  $\mathcal{L}_{+}$  as follows. First, we note that the tangent space of  $\mathcal{L}_{+}$  at a given  $L \in \mathcal{L}_{+}$  is identified with the linear space  $\mathcal{L}$ . For such tangent space, we treat the strict lower triangular space  $\lfloor \mathcal{L} \rfloor \coloneqq \{\lfloor X \rfloor : X \in \mathcal{L}\}$  as the Euclidean space  $\mathbb{R}^{m(m-1)/2}$  with the usual Frobenius inner product  $\langle X, Y \rangle_{\mathrm{F}} = \sum_{i,j=1}^{m} X_{ij} Y_{ij}$  for all  $X, Y \in \lfloor \mathcal{L} \rfloor$ . For the diagonal part  $\mathbb{D}(\mathcal{L}) \coloneqq \{\mathbb{D}(Z) : Z \in \mathcal{L}\}$ , we equipped it with a different inner product defined by  $\langle \mathbb{D}(L)^{-1} \mathbb{D}(X), \mathbb{D}(L)^{-1} \mathbb{D}(Y) \rangle_{\mathrm{F}}$ . Finally, combining these two components together, we define a metric  $\tilde{g}$  for tangent spaces  $T_L \mathcal{L}_{+}$  (identified with  $\mathcal{L}$ ) by

$$
\begin{array}{l} \tilde {g} _ {L} (X, Y) = \langle \lfloor X \rfloor , \lfloor Y \rfloor \rangle_ {\mathrm {F}} + \langle \mathbb {D} (L) ^ {- 1} \mathbb {D} (X), \mathbb {D} (L) ^ {- 1} \mathbb {D} (Y) \rangle_ {\mathrm {F}} \\ = \sum_ {i > j} X _ {i j} Y _ {i j} + \sum_ {j = 1} ^ {m} X _ {j j} Y _ {j j} L _ {j j} ^ {- 2}. \\ \end{array}
$$

It is straightforward to show that the space  $\mathcal{L}_{+}$ , equipped with the metric  $\tilde{g}$ , is a Riemannian manifold. We begin with geodesics on the manifold  $(\mathcal{L}_{+},\tilde{g})$  to investigate its basic properties.

$$
d _ {\mathcal {L} _ {+}} (L, K) = \{\tilde {g} _ {L} (X, X) \} ^ {1 / 2} = \left\{\sum_ {i > j} (L _ {i j} - K _ {i j}) ^ {2} + \sum_ {j = 1} ^ {m} (\log L _ {j j} - \log K _ {j j}) ^ {2} \right\} ^ {1 / 2},
$$

where  $X$  is the same as the above. The expression for the distance function can be equivalently and more compactly written as

$$
d _ {\mathcal {L} _ {+}} (L, K) = \{\| \lfloor L \rfloor - \lfloor K \rfloor \| _ {\mathrm {F}} ^ {2} + \| \log \mathbb {D} (L) - \log \mathbb {D} (K) \| _ {\mathrm {F}} ^ {2} \} ^ {1 / 2}.
$$

$$
\gamma_ {P, W} (t) = \mathcal {S} (\tilde {\gamma} _ {L, X} (t)) = \tilde {\gamma} _ {L, X} (t) \tilde {\gamma} _ {L, X} ^ {\top} (t),
$$

where  $X = L(L^{-1}WL^{-\top})_{\frac{1}{2}}$  and  $W \in T_P\mathcal{S}_m^+$ . Similarly, the Riemannian exponential at  $P$  is given by

$$
\mathrm {E x p} _ {P} W = \mathcal {S} (\widetilde {\mathrm {E x p}} _ {L} X) = (\widetilde {\mathrm {E x p}} _ {L} X) (\widetilde {\mathrm {E x p}} _ {L} X) ^ {\top},
$$

where  $X = L(L^{-1}WL^{-\top})_{\frac{1}{2}}$  and  $W \in T_P\mathcal{S}_m^+$ . Similarly, the Riemannian exponential at  $P$  is given by

$$
\mathrm {E x p} _ {P} W = \mathcal {S} (\widetilde {\mathrm {E x p}} _ {L} X) = (\widetilde {\mathrm {E x p}} _ {L} X) (\widetilde {\mathrm {E x p}} _ {L} X) ^ {\top},
$$

while the geodesic between  $P$  and  $Q$  is characterized by

inverse of  $L$ , denoted by  $L_{\odot}^{-1}$ , is  $\mathbb{D}(L)^{-1} - \lfloor L\rfloor$ . The left translation by  $A\in \mathcal{L}_{+}$  is denoted by  $\ell_A:B\mapsto A\odot B$ . One can check that the differential of this operation at  $L\in \mathcal{L}_{+}$  is

$$
D _ {L} \ell_ {A}: X \mapsto \lfloor X \rfloor + \mathbb {D} (A) \mathbb {D} (X), \tag {3.2}
$$

where it is noted that the differential  $D_{L}\ell_{A}$  does not depend on  $L$ . Given the above expression, one can find that

$$
D _ {L} \ell_ {A}: X \mapsto \lfloor X \rfloor + \mathbb {D} (A) \mathbb {D} (X), \tag {3.2}
$$

where it is noted that the differential  $D_{L}\ell_{A}$  does not depend on  $L$ . Given the above expression, one can find that

$$
\begin{array}{l} \tilde {g} _ {A \odot L} ((D _ {L} \ell_ {A}) (X), (D _ {L} \ell_ {A}) (Y)) \\ = \tilde {g} _ {A \odot L} (\left\lfloor X \right\rfloor + \mathbb {D} (A) \mathbb {D} (X), \left\lfloor Y \right\rfloor + \mathbb {D} (A) \mathbb {D} (Y)) \\ = \tilde {g} _ {L} (X, Y). \\ \end{array}
$$

The group operator  $\odot$  and maps  $\mathcal{S}$  and  $\mathcal{L}$  together induce a smooth operation  $\otimes$  on  $S_{m}^{+}$ , defined by

$$
\begin{array}{l} P \circledast Q = \mathcal {S} (\mathcal {L} (P) \circledast \mathcal {L} (Q)) \\ = \left(\mathcal {L} (P) \circledast \mathcal {L} (Q)\right) \left(\mathcal {L} (P) \circledast \mathcal {L} (Q)\right) ^ {\top} \quad \text {f o r} P, Q \in S _ {m} ^ {+}. \\ \end{array}
$$

In addition, both  $\mathcal{L}$  and  $\mathcal{S}$  are Riemannian isometries and group isomorphisms between Lie groups  $(\mathcal{L}_{+},\tilde{g},\circledast)$  and  $(\mathcal{S}_m^+,g,\bullet)$ .

$$
\begin{array}{l} P \circledast Q = \mathcal {S} (\mathcal {L} (P) \circledast \mathcal {L} (Q)) \\ = \left(\mathcal {L} (P) \circledast \mathcal {L} (Q)\right) \left(\mathcal {L} (P) \circledast \mathcal {L} (Q)\right) ^ {\top} \quad \text {f o r} P, Q \in S _ {m} ^ {+}. \\ \end{array}
$$

In addition, both  $\mathcal{L}$  and  $\mathcal{S}$  are Riemannian isometries and group isomorphisms between Lie groups  $(\mathcal{L}_{+},\tilde{g},\circledast)$  and  $(\mathcal{S}_m^+,g,\bullet)$ .

THEOREM 5. The space  $(S_m^+, \otimes)$  is an abelian Lie group. Moreover, the metric  $g$  defined in (3.1) is a bi-invariant metric on  $(S_m^+, \otimes)$ .

$$
\begin{array}{l} Z (h q) = \left(D _ {p} \ell_ {h q p ^ {- 1}}\right) u = \left\{D _ {p} \left(\ell_ {h} \circ \ell_ {q p ^ {- 1}}\right) \right\} u \\ = \left(D _ {q} \ell_ {h}\right) \left(D _ {p} \ell_ {q p ^ {- 1}}\right) u = \left(D _ {q} \ell_ {h}\right) (Z (q)), \\ \end{array}
$$

where the third equality is obtained by the chain rule of differential. Consequently,  $\nabla_{\gamma_p'}Z = 0$  since  $\gamma_p'(t) = Y(p\gamma_e(t))$  and  $\nabla_Y Z = 0$  for  $Y,Z\in \mathfrak{g}$ . As additionally  $Z(\gamma_p(0)) = Z(p) = u$ , transportation of  $u$  along the geodesic  $\gamma_{p}$  is realized by the left-invariant vector field  $Z$ . Since  $\gamma_{p}$  is a geodesic with  $\gamma_p(0) = p$  and  $\gamma_p(1) = \ell_p\mathfrak{e}\mathfrak{p}(Y) = p(p^{-1}q) = q$ , we have that

$$
\tau_ {p, q} (u) = Z \left(\gamma_ {p} (1)\right) = Z (q) = \left(D _ {p} \ell_ {q p ^ {- 1}}\right) u,
$$

where the third equality is obtained by the chain rule of differential. Consequently,  $\nabla_{\gamma_p'}Z = 0$  since  $\gamma_p'(t) = Y(p\gamma_e(t))$  and  $\nabla_Y Z = 0$  for  $Y,Z\in \mathfrak{g}$ . As additionally  $Z(\gamma_p(0)) = Z(p) = u$ , transportation of  $u$  along the geodesic  $\gamma_{p}$  is realized by the left-invariant vector field  $Z$ . Since  $\gamma_{p}$  is a geodesic with  $\gamma_p(0) = p$  and  $\gamma_p(1) = \ell_p\mathfrak{e}\mathfrak{p}(Y) = p(p^{-1}q) = q$ , we have that

$$
\tau_ {p, q} (u) = Z \left(\gamma_ {p} (1)\right) = Z (q) = \left(D _ {p} \ell_ {q p ^ {- 1}}\right) u,
$$

as claimed.

$$
\begin{array}{l} \det  \mathbb {E} L = \exp [ \operatorname {t r} \{\log \mathbb {D} (\mathbb {E} L) \} ] = \exp [ \operatorname {t r} \{\mathbb {E} \log \mathbb {D} (L) \} ] \\ = \exp \left[ \mathbb {E} \operatorname {t r} \left\{\log \mathbb {D} (L) \right\} \right] = \exp \left\{\mathbb {E} \log (\det  L) \right\}, \\ \end{array}
$$

which proves the case of  $(\mathcal{L}_{+},\tilde{g})$

For the case that  $Q$  is a random element on  $(\mathcal{S}_m^+, g)$ , we denote it by  $P$  and observe that  $\operatorname{det}(AB) = (\operatorname{det} A)(\operatorname{det} B)$  for any square matrices  $A$  and  $B$ . Let  $L = \mathcal{L}(P)$ . Then, one has

which proves the case of  $(\mathcal{L}_{+},\tilde{g})$

For the case that  $Q$  is a random element on  $(\mathcal{S}_m^+, g)$ , we denote it by  $P$  and observe that  $\operatorname{det}(AB) = (\operatorname{det} A)(\operatorname{det} B)$  for any square matrices  $A$  and  $B$ . Let  $L = \mathcal{L}(P)$ . Then, one has

$$
\begin{array}{l} \det  \mathbb {E} P = \det  \left\{\left(\mathbb {E} L\right) \left(\mathbb {E} L\right) ^ {\top} \right\} \\ = \left\{\det  (\mathbb {E} L) \right\} \left\{\det  (\mathbb {E} L) ^ {\top} \right\} \\ = \exp \left\{\mathbb {E} \log (\det  L) \right\} \exp \left\{\mathbb {E} \log (\det  L ^ {\top}) \right\} \\ = \exp \left\{\mathbb {E} \left(\log \det  L + \log \det  L ^ {\top}\right) \right\} \\ = \exp \left[ \mathbb {E} \log \left\{\left(\det  L\right) \det  \left(L ^ {\top}\right) \right\} \right], \\ \end{array}
$$

For the case that  $Q$  is a random element on  $(\mathcal{S}_m^+, g)$ , we denote it by  $P$  and observe that  $\operatorname{det}(AB) = (\operatorname{det} A)(\operatorname{det} B)$  for any square matrices  $A$  and  $B$ . Let  $L = \mathcal{L}(P)$ . Then, one has

$$
\begin{array}{l} \det  \mathbb {E} P = \det  \left\{\left(\mathbb {E} L\right) \left(\mathbb {E} L\right) ^ {\top} \right\} \\ = \left\{\det  (\mathbb {E} L) \right\} \left\{\det  (\mathbb {E} L) ^ {\top} \right\} \\ = \exp \left\{\mathbb {E} \log (\det  L) \right\} \exp \left\{\mathbb {E} \log (\det  L ^ {\top}) \right\} \\ = \exp \left\{\mathbb {E} \left(\log \det  L + \log \det  L ^ {\top}\right) \right\} \\ = \exp \left[ \mathbb {E} \log \left\{\left(\det  L\right) \det  \left(L ^ {\top}\right) \right\} \right], \\ \end{array}
$$

$$
\begin{array}{l} = \exp \left\{\mathbb {E} \log \det  (L L ^ {\top}) \right\} \\ = \exp \left\{\mathbb {E} \log \det  (P) \right\}, \\ \end{array}
$$


\begin{proposition}[Proposition 8]
PROPOSITION 8. The sectional curvature of  $(\mathcal{L}_{+},\tilde{g})$  and  $(S_m^+,g)$  is constantly zero.
\end{proposition}

\begin{proof}
Proof. For  $(\mathcal{L}_{+},\tilde{g})$ , in the proof of Proposition 3, it has been shown that all Christoffel symbols are zero under the selected coordinate. This implies that the Riemannian curvature tensor is identically zero and hence so is the sectional curvature. The case of  $(S_m^+,g)$  follows from the fact that  $\mathcal{S}$  is an isometry that preserves sectional curvature.

PROPOSITION 9. If  $L$  is a random element on  $(\mathcal{L}_{+},\tilde{g})$  and  $\mathbb{E}d_{\mathcal{L}_+}^2 (A,L) < \infty$  for some  $A\in \mathcal{L}_{+}$ . Then the Fréchet mean  $\mathbb{E}L$  exists and is unique. Similarly, the Log-Cholesky mean of a random SPD matrix  $P$  exists and is unique if  $\mathbb{E}d_{S_m^+}^2 (S,P) < \infty$  for some SPD matrix  $S$ .

\end{proof}

\paragraph{Background.}
-  $\exp(S)$  is an SPD matrix for a symmetric matrix  $S$ , while  $\log(P)$  is a symmetric metric for an SPD matrix  $P$ .

- Diagonal elements of an SPD matrix are all positive. This can be seen from the fact that  $P_{jj} = e_j^\top Pe_j > 0$  for  $P \in S_m^+$ , where  $e_j$  is the unit vector with 1 at the  $j$ th coordinate and 0 elsewhere.

2.2. Cholesky decomposition. Cholesky decomposition, named after André-Louis Cholesky, represents a real  $m \times m$  SPD matrix  $P$  as a product of a lower triangular matrix  $L$  and its transpose, i.e.,  $P = LL^{\top}$ . If the diagonal elements of  $L$  are restricted to be positive, then the decomposition is unique according to Theorem 4.2.5 of [16]. Such lower triangular matrix, denoted by  $\mathcal{L}(P)$ , is called the Cholesky factor of  $P$ . Since in addition  $L = \mathcal{L}(LL^{\top})$  for each  $L \in \mathcal{L}_{+}$ , the map  $\mathcal{L}: S_{m}^{+} \to \mathcal{L}_{+}$  is bijective. In other words, there is one-to-one correspondence between SPD matrices and lower triangular matrices whose diagonal elements are all positive.

3.1. Riemannian geometry on Cholesky spaces. For matrices in  $\mathcal{L}_{+}$ , as off-diagonal elements in the lower triangular part are unconstrained while diagonal ones are restricted to be positive,  $\mathcal{L}_{+}$  can be parameterized by  $\mathcal{L} \ni X \coloneqq \varphi(X) \in \mathcal{L}_{+}$  in the way that  $(\varphi(X))_{ij} = X_{ij}$  if  $i \neq j$  and  $(\varphi(X))_{jj} = \exp(X_{jj})$ . This motivates us to, respectively, endow the unconstrained part  $\lfloor \cdot \rfloor$  and the positive diagonal part  $\mathbb{D}(\cdot)$  with a different metric and then combine them into a Riemannian metric on  $\mathcal{L}_{+}$  as follows. First, we note that the tangent space of  $\mathcal{L}_{+}$  at a given  $L \in \mathcal{L}_{+}$  is identified with the linear space  $\mathcal{L}$ . For such tangent space, we treat the strict lower triangular space  $\lfloor \mathcal{L} \rfloor \coloneqq \{\lfloor X \rfloor : X \in \mathcal{L}\}$  as the Euclidean space  $\mathbb{R}^{m(m-1)/2}$  with the usual Frobenius inner product  $\langle X, Y \rangle_{\mathrm{F}} = \sum_{i,j=1}^{m} X_{ij} Y_{ij}$  for all  $X, Y \in \lfloor \mathcal{L} \rfloor$ . For the diagonal part  $\mathbb{D}(\mathcal{L}) \coloneqq \{\mathbb{D}(Z) : Z \in \mathcal{L}\}$ , we equipped it with a different inner product defined by  $\langle \mathbb{D}(L)^{-1} \mathbb{D}(X), \mathbb{D}(L)^{-1} \mathbb{D}(Y) \rangle_{\mathrm{F}}$ . Finally, combining these two components together, we define a metric  $\tilde{g}$  for tangent spaces  $T_L \mathcal{L}_{+}$  (identified with  $\mathcal{L}$ ) by

$$
\begin{array}{l} \tilde {g} _ {L} (X, Y) = \langle \lfloor X \rfloor , \lfloor Y \rfloor \rangle_ {\mathrm {F}} + \langle \mathbb {D} (L) ^ {- 1} \mathbb {D} (X), \mathbb {D} (L) ^ {- 1} \mathbb {D} (Y) \rangle_ {\mathrm {F}} \\ = \sum_ {i > j} X _ {i j} Y _ {i j} + \sum_ {j = 1} ^ {m} X _ {j j} Y _ {j j} L _ {j j} ^ {- 2}. \\ \end{array}
$$

It is straightforward to show that the space  $\mathcal{L}_{+}$ , equipped with the metric  $\tilde{g}$ , is a Riemannian manifold. We begin with geodesics on the manifold  $(\mathcal{L}_{+},\tilde{g})$  to investigate its basic properties.

$$
L ^ {- 1} W L ^ {- \top} = X ^ {\top} L ^ {- \top} + L ^ {- 1} X = L ^ {- 1} X + (L ^ {- 1} X) ^ {\top}.
$$

Note that  $L^{-1}X$  is also a lower triangular matrix, and the matrix on the left-hand side is symmetric, we deduce that  $(L^{-1}WL^{-\top})_{\frac{1}{2}} = L^{-1}X$ , which gives rise to  $X = L(L^{-1}WL^{-\top})_{\frac{1}{2}}$ . The linear map  $(D_L\mathcal{S})(X) = LX^\top +XL^\top$  is one-to-one, since from the above derivation,  $(D_L\mathcal{S})(X) = 0$  if and only if  $X = 0$ .

Given the above proposition, the manifold map  $\mathcal{S}$ , which is exactly the inverse of the Cholesky map  $\mathcal{L}$  discussed in subsection 2.2, induces a Riemannian metric on  $S_{m}^{+}$ , denoted by  $g$  and called Log-Cholesky metric, given by

Note that  $L^{-1}X$  is also a lower triangular matrix, and the matrix on the left-hand side is symmetric, we deduce that  $(L^{-1}WL^{-\top})_{\frac{1}{2}} = L^{-1}X$ , which gives rise to  $X = L(L^{-1}WL^{-\top})_{\frac{1}{2}}$ . The linear map  $(D_L\mathcal{S})(X) = LX^\top +XL^\top$  is one-to-one, since from the above derivation,  $(D_L\mathcal{S})(X) = 0$  if and only if  $X = 0$ .

Given the above proposition, the manifold map  $\mathcal{S}$ , which is exactly the inverse of the Cholesky map  $\mathcal{L}$  discussed in subsection 2.2, induces a Riemannian metric on  $S_{m}^{+}$ , denoted by  $g$  and called Log-Cholesky metric, given by

$$
g _ {P} (W, V) = \tilde {g} _ {L} \left(L \left(L ^ {- 1} W L ^ {- \top}\right) _ {\frac {1}{2}}, L \left(L ^ {- 1} V L ^ {- \top}\right) _ {\frac {1}{2}}\right), \tag {3.1}
$$

inverse of  $L$ , denoted by  $L_{\odot}^{-1}$ , is  $\mathbb{D}(L)^{-1} - \lfloor L\rfloor$ . The left translation by  $A\in \mathcal{L}_{+}$  is denoted by  $\ell_A:B\mapsto A\odot B$ . One can check that the differential of this operation at  $L\in \mathcal{L}_{+}$  is

$$
D _ {L} \ell_ {A}: X \mapsto \lfloor X \rfloor + \mathbb {D} (A) \mathbb {D} (X), \tag {3.2}
$$

where it is noted that the differential  $D_{L}\ell_{A}$  does not depend on  $L$ . Given the above expression, one can find that

$$
D _ {L} \ell_ {A}: X \mapsto \lfloor X \rfloor + \mathbb {D} (A) \mathbb {D} (X), \tag {3.2}
$$

where it is noted that the differential  $D_{L}\ell_{A}$  does not depend on  $L$ . Given the above expression, one can find that

$$
\begin{array}{l} \tilde {g} _ {A \odot L} ((D _ {L} \ell_ {A}) (X), (D _ {L} \ell_ {A}) (Y)) \\ = \tilde {g} _ {A \odot L} (\left\lfloor X \right\rfloor + \mathbb {D} (A) \mathbb {D} (X), \left\lfloor Y \right\rfloor + \mathbb {D} (A) \mathbb {D} (Y)) \\ = \tilde {g} _ {L} (X, Y). \\ \end{array}
$$

$$
\begin{array}{l} P \circledast Q = \mathcal {S} (\mathcal {L} (P) \circledast \mathcal {L} (Q)) \\ = \left(\mathcal {L} (P) \circledast \mathcal {L} (Q)\right) \left(\mathcal {L} (P) \circledast \mathcal {L} (Q)\right) ^ {\top} \quad \text {f o r} P, Q \in S _ {m} ^ {+}. \\ \end{array}
$$

In addition, both  $\mathcal{L}$  and  $\mathcal{S}$  are Riemannian isometries and group isomorphisms between Lie groups  $(\mathcal{L}_{+},\tilde{g},\circledast)$  and  $(\mathcal{S}_m^+,g,\bullet)$ .

THEOREM 5. The space  $(S_m^+, \otimes)$  is an abelian Lie group. Moreover, the metric  $g$  defined in (3.1) is a bi-invariant metric on  $(S_m^+, \otimes)$ .

$$
\begin{array}{l} \det  \mathbb {E} L = \exp [ \operatorname {t r} \{\log \mathbb {D} (\mathbb {E} L) \} ] = \exp [ \operatorname {t r} \{\mathbb {E} \log \mathbb {D} (L) \} ] \\ = \exp \left[ \mathbb {E} \operatorname {t r} \left\{\log \mathbb {D} (L) \right\} \right] = \exp \left\{\mathbb {E} \log (\det  L) \right\}, \\ \end{array}
$$

which proves the case of  $(\mathcal{L}_{+},\tilde{g})$

For the case that  $Q$  is a random element on  $(\mathcal{S}_m^+, g)$ , we denote it by  $P$  and observe that  $\operatorname{det}(AB) = (\operatorname{det} A)(\operatorname{det} B)$  for any square matrices  $A$  and  $B$ . Let  $L = \mathcal{L}(P)$ . Then, one has

which proves the case of  $(\mathcal{L}_{+},\tilde{g})$

For the case that  $Q$  is a random element on  $(\mathcal{S}_m^+, g)$ , we denote it by  $P$  and observe that  $\operatorname{det}(AB) = (\operatorname{det} A)(\operatorname{det} B)$  for any square matrices  $A$  and  $B$ . Let  $L = \mathcal{L}(P)$ . Then, one has

$$
\begin{array}{l} \det  \mathbb {E} P = \det  \left\{\left(\mathbb {E} L\right) \left(\mathbb {E} L\right) ^ {\top} \right\} \\ = \left\{\det  (\mathbb {E} L) \right\} \left\{\det  (\mathbb {E} L) ^ {\top} \right\} \\ = \exp \left\{\mathbb {E} \log (\det  L) \right\} \exp \left\{\mathbb {E} \log (\det  L ^ {\top}) \right\} \\ = \exp \left\{\mathbb {E} \left(\log \det  L + \log \det  L ^ {\top}\right) \right\} \\ = \exp \left[ \mathbb {E} \log \left\{\left(\det  L\right) \det  \left(L ^ {\top}\right) \right\} \right], \\ \end{array}
$$


\begin{proposition}[Proposition 9]
PROPOSITION 9. If  $L$  is a random element on  $(\mathcal{L}_{+},\tilde{g})$  and  $\mathbb{E}d_{\mathcal{L}_+}^2 (A,L) < \infty$  for some  $A\in \mathcal{L}_{+}$ . Then the Fréchet mean  $\mathbb{E}L$  exists and is unique. Similarly, the Log-Cholesky mean of a random SPD matrix  $P$  exists and is unique if  $\mathbb{E}d_{S_m^+}^2 (S,P) < \infty$  for some SPD matrix  $S$ .
\end{proposition}

\begin{proof}
Proof. In the case of  $(\mathcal{L}_{+},\tilde{g})$ , we define  $\psi (L) = \lfloor L\rfloor +\log \mathbb{D}(L)$ . It can be shown that  $\psi$  is a diffeomorphism (and hence also a homeomorphism) between  $\mathcal{L}_{+}$  and  $\mathcal{L}$ . Therefore,  $\mathcal{L}_{+}$  is simply connected since  $\mathcal{L}$  is. In Proposition 8 we have shown that  $\mathcal{L}_{+}$  has zero sectional curvature. Thus, the existence and uniqueness of Fréchet mean follows from Theorem 2.1 of [8]. The case of  $(S_m^+,g)$  follows from the isometry of  $\mathcal{S}$ .  $\square$

The next result shows that the Log-Cholesky mean of a random SPD matrix is computable from the Fréchet mean of its Cholesky factor. Also, it characterizes the Fréchet mean of a random Cholesky factor, which is important for us to derive a closed form for the Log-Cholesky average of a finite number of matrices in the next subsection.

\end{proof}

\paragraph{Background.}
- For any invertible matrix  $X$ , the matrix  $XX^{\top}$  is an SPD matrix.

-  $\exp(S)$  is an SPD matrix for a symmetric matrix  $S$ , while  $\log(P)$  is a symmetric metric for an SPD matrix  $P$ .

- Diagonal elements of an SPD matrix are all positive. This can be seen from the fact that  $P_{jj} = e_j^\top Pe_j > 0$  for  $P \in S_m^+$ , where  $e_j$  is the unit vector with 1 at the  $j$ th coordinate and 0 elsewhere.

-  $\exp(S)$  is an SPD matrix for a symmetric matrix  $S$ , while  $\log(P)$  is a symmetric metric for an SPD matrix  $P$ .

- Diagonal elements of an SPD matrix are all positive. This can be seen from the fact that  $P_{jj} = e_j^\top Pe_j > 0$  for  $P \in S_m^+$ , where  $e_j$  is the unit vector with 1 at the  $j$ th coordinate and 0 elsewhere.

2.2. Cholesky decomposition. Cholesky decomposition, named after André-Louis Cholesky, represents a real  $m \times m$  SPD matrix  $P$  as a product of a lower triangular matrix  $L$  and its transpose, i.e.,  $P = LL^{\top}$ . If the diagonal elements of  $L$  are restricted to be positive, then the decomposition is unique according to Theorem 4.2.5 of [16]. Such lower triangular matrix, denoted by  $\mathcal{L}(P)$ , is called the Cholesky factor of  $P$ . Since in addition  $L = \mathcal{L}(LL^{\top})$  for each  $L \in \mathcal{L}_{+}$ , the map  $\mathcal{L}: S_{m}^{+} \to \mathcal{L}_{+}$  is bijective. In other words, there is one-to-one correspondence between SPD matrices and lower triangular matrices whose diagonal elements are all positive.

3.1. Riemannian geometry on Cholesky spaces. For matrices in  $\mathcal{L}_{+}$ , as off-diagonal elements in the lower triangular part are unconstrained while diagonal ones are restricted to be positive,  $\mathcal{L}_{+}$  can be parameterized by  $\mathcal{L} \ni X \coloneqq \varphi(X) \in \mathcal{L}_{+}$  in the way that  $(\varphi(X))_{ij} = X_{ij}$  if  $i \neq j$  and  $(\varphi(X))_{jj} = \exp(X_{jj})$ . This motivates us to, respectively, endow the unconstrained part  $\lfloor \cdot \rfloor$  and the positive diagonal part  $\mathbb{D}(\cdot)$  with a different metric and then combine them into a Riemannian metric on  $\mathcal{L}_{+}$  as follows. First, we note that the tangent space of  $\mathcal{L}_{+}$  at a given  $L \in \mathcal{L}_{+}$  is identified with the linear space  $\mathcal{L}$ . For such tangent space, we treat the strict lower triangular space  $\lfloor \mathcal{L} \rfloor \coloneqq \{\lfloor X \rfloor : X \in \mathcal{L}\}$  as the Euclidean space  $\mathbb{R}^{m(m-1)/2}$  with the usual Frobenius inner product  $\langle X, Y \rangle_{\mathrm{F}} = \sum_{i,j=1}^{m} X_{ij} Y_{ij}$  for all  $X, Y \in \lfloor \mathcal{L} \rfloor$ . For the diagonal part  $\mathbb{D}(\mathcal{L}) \coloneqq \{\mathbb{D}(Z) : Z \in \mathcal{L}\}$ , we equipped it with a different inner product defined by  $\langle \mathbb{D}(L)^{-1} \mathbb{D}(X), \mathbb{D}(L)^{-1} \mathbb{D}(Y) \rangle_{\mathrm{F}}$ . Finally, combining these two components together, we define a metric  $\tilde{g}$  for tangent spaces  $T_L \mathcal{L}_{+}$  (identified with  $\mathcal{L}$ ) by

$$
\begin{array}{l} \tilde {g} _ {L} (X, Y) = \langle \lfloor X \rfloor , \lfloor Y \rfloor \rangle_ {\mathrm {F}} + \langle \mathbb {D} (L) ^ {- 1} \mathbb {D} (X), \mathbb {D} (L) ^ {- 1} \mathbb {D} (Y) \rangle_ {\mathrm {F}} \\ = \sum_ {i > j} X _ {i j} Y _ {i j} + \sum_ {j = 1} ^ {m} X _ {j j} Y _ {j j} L _ {j j} ^ {- 2}. \\ \end{array}
$$

It is straightforward to show that the space  $\mathcal{L}_{+}$ , equipped with the metric  $\tilde{g}$ , is a Riemannian manifold. We begin with geodesics on the manifold  $(\mathcal{L}_{+},\tilde{g})$  to investigate its basic properties.

$$
L ^ {- 1} W L ^ {- \top} = X ^ {\top} L ^ {- \top} + L ^ {- 1} X = L ^ {- 1} X + (L ^ {- 1} X) ^ {\top}.
$$

Note that  $L^{-1}X$  is also a lower triangular matrix, and the matrix on the left-hand side is symmetric, we deduce that  $(L^{-1}WL^{-\top})_{\frac{1}{2}} = L^{-1}X$ , which gives rise to  $X = L(L^{-1}WL^{-\top})_{\frac{1}{2}}$ . The linear map  $(D_L\mathcal{S})(X) = LX^\top +XL^\top$  is one-to-one, since from the above derivation,  $(D_L\mathcal{S})(X) = 0$  if and only if  $X = 0$ .

Given the above proposition, the manifold map  $\mathcal{S}$ , which is exactly the inverse of the Cholesky map  $\mathcal{L}$  discussed in subsection 2.2, induces a Riemannian metric on  $S_{m}^{+}$ , denoted by  $g$  and called Log-Cholesky metric, given by

Note that  $L^{-1}X$  is also a lower triangular matrix, and the matrix on the left-hand side is symmetric, we deduce that  $(L^{-1}WL^{-\top})_{\frac{1}{2}} = L^{-1}X$ , which gives rise to  $X = L(L^{-1}WL^{-\top})_{\frac{1}{2}}$ . The linear map  $(D_L\mathcal{S})(X) = LX^\top +XL^\top$  is one-to-one, since from the above derivation,  $(D_L\mathcal{S})(X) = 0$  if and only if  $X = 0$ .

Given the above proposition, the manifold map  $\mathcal{S}$ , which is exactly the inverse of the Cholesky map  $\mathcal{L}$  discussed in subsection 2.2, induces a Riemannian metric on  $S_{m}^{+}$ , denoted by  $g$  and called Log-Cholesky metric, given by

$$
g _ {P} (W, V) = \tilde {g} _ {L} \left(L \left(L ^ {- 1} W L ^ {- \top}\right) _ {\frac {1}{2}}, L \left(L ^ {- 1} V L ^ {- \top}\right) _ {\frac {1}{2}}\right), \tag {3.1}
$$

$$
\gamma_ {P, W} (t) = \mathcal {S} (\tilde {\gamma} _ {L, X} (t)) = \tilde {\gamma} _ {L, X} (t) \tilde {\gamma} _ {L, X} ^ {\top} (t),
$$

where  $X = L(L^{-1}WL^{-\top})_{\frac{1}{2}}$  and  $W \in T_P\mathcal{S}_m^+$ . Similarly, the Riemannian exponential at  $P$  is given by

$$
\mathrm {E x p} _ {P} W = \mathcal {S} (\widetilde {\mathrm {E x p}} _ {L} X) = (\widetilde {\mathrm {E x p}} _ {L} X) (\widetilde {\mathrm {E x p}} _ {L} X) ^ {\top},
$$

where  $X = L(L^{-1}WL^{-\top})_{\frac{1}{2}}$  and  $W \in T_P\mathcal{S}_m^+$ . Similarly, the Riemannian exponential at  $P$  is given by

$$
\mathrm {E x p} _ {P} W = \mathcal {S} (\widetilde {\mathrm {E x p}} _ {L} X) = (\widetilde {\mathrm {E x p}} _ {L} X) (\widetilde {\mathrm {E x p}} _ {L} X) ^ {\top},
$$

while the geodesic between  $P$  and  $Q$  is characterized by

inverse of  $L$ , denoted by  $L_{\odot}^{-1}$ , is  $\mathbb{D}(L)^{-1} - \lfloor L\rfloor$ . The left translation by  $A\in \mathcal{L}_{+}$  is denoted by  $\ell_A:B\mapsto A\odot B$ . One can check that the differential of this operation at  $L\in \mathcal{L}_{+}$  is

$$
D _ {L} \ell_ {A}: X \mapsto \lfloor X \rfloor + \mathbb {D} (A) \mathbb {D} (X), \tag {3.2}
$$

where it is noted that the differential  $D_{L}\ell_{A}$  does not depend on  $L$ . Given the above expression, one can find that

$$
D _ {L} \ell_ {A}: X \mapsto \lfloor X \rfloor + \mathbb {D} (A) \mathbb {D} (X), \tag {3.2}
$$

where it is noted that the differential  $D_{L}\ell_{A}$  does not depend on  $L$ . Given the above expression, one can find that

$$
\begin{array}{l} \tilde {g} _ {A \odot L} ((D _ {L} \ell_ {A}) (X), (D _ {L} \ell_ {A}) (Y)) \\ = \tilde {g} _ {A \odot L} (\left\lfloor X \right\rfloor + \mathbb {D} (A) \mathbb {D} (X), \left\lfloor Y \right\rfloor + \mathbb {D} (A) \mathbb {D} (Y)) \\ = \tilde {g} _ {L} (X, Y). \\ \end{array}
$$

$$
\begin{array}{l} P \circledast Q = \mathcal {S} (\mathcal {L} (P) \circledast \mathcal {L} (Q)) \\ = \left(\mathcal {L} (P) \circledast \mathcal {L} (Q)\right) \left(\mathcal {L} (P) \circledast \mathcal {L} (Q)\right) ^ {\top} \quad \text {f o r} P, Q \in S _ {m} ^ {+}. \\ \end{array}
$$

In addition, both  $\mathcal{L}$  and  $\mathcal{S}$  are Riemannian isometries and group isomorphisms between Lie groups  $(\mathcal{L}_{+},\tilde{g},\circledast)$  and  $(\mathcal{S}_m^+,g,\bullet)$ .

THEOREM 5. The space  $(S_m^+, \otimes)$  is an abelian Lie group. Moreover, the metric  $g$  defined in (3.1) is a bi-invariant metric on  $(S_m^+, \otimes)$ .

$$
\begin{array}{l} \det  \mathbb {E} L = \exp [ \operatorname {t r} \{\log \mathbb {D} (\mathbb {E} L) \} ] = \exp [ \operatorname {t r} \{\mathbb {E} \log \mathbb {D} (L) \} ] \\ = \exp \left[ \mathbb {E} \operatorname {t r} \left\{\log \mathbb {D} (L) \right\} \right] = \exp \left\{\mathbb {E} \log (\det  L) \right\}, \\ \end{array}
$$

which proves the case of  $(\mathcal{L}_{+},\tilde{g})$

For the case that  $Q$  is a random element on  $(\mathcal{S}_m^+, g)$ , we denote it by  $P$  and observe that  $\operatorname{det}(AB) = (\operatorname{det} A)(\operatorname{det} B)$  for any square matrices  $A$  and  $B$ . Let  $L = \mathcal{L}(P)$ . Then, one has

which proves the case of  $(\mathcal{L}_{+},\tilde{g})$

For the case that  $Q$  is a random element on  $(\mathcal{S}_m^+, g)$ , we denote it by  $P$  and observe that  $\operatorname{det}(AB) = (\operatorname{det} A)(\operatorname{det} B)$  for any square matrices  $A$  and  $B$ . Let  $L = \mathcal{L}(P)$ . Then, one has

$$
\begin{array}{l} \det  \mathbb {E} P = \det  \left\{\left(\mathbb {E} L\right) \left(\mathbb {E} L\right) ^ {\top} \right\} \\ = \left\{\det  (\mathbb {E} L) \right\} \left\{\det  (\mathbb {E} L) ^ {\top} \right\} \\ = \exp \left\{\mathbb {E} \log (\det  L) \right\} \exp \left\{\mathbb {E} \log (\det  L ^ {\top}) \right\} \\ = \exp \left\{\mathbb {E} \left(\log \det  L + \log \det  L ^ {\top}\right) \right\} \\ = \exp \left[ \mathbb {E} \log \left\{\left(\det  L\right) \det  \left(L ^ {\top}\right) \right\} \right], \\ \end{array}
$$

For the case that  $Q$  is a random element on  $(\mathcal{S}_m^+, g)$ , we denote it by  $P$  and observe that  $\operatorname{det}(AB) = (\operatorname{det} A)(\operatorname{det} B)$  for any square matrices  $A$  and  $B$ . Let  $L = \mathcal{L}(P)$ . Then, one has

$$
\begin{array}{l} \det  \mathbb {E} P = \det  \left\{\left(\mathbb {E} L\right) \left(\mathbb {E} L\right) ^ {\top} \right\} \\ = \left\{\det  (\mathbb {E} L) \right\} \left\{\det  (\mathbb {E} L) ^ {\top} \right\} \\ = \exp \left\{\mathbb {E} \log (\det  L) \right\} \exp \left\{\mathbb {E} \log (\det  L ^ {\top}) \right\} \\ = \exp \left\{\mathbb {E} \left(\log \det  L + \log \det  L ^ {\top}\right) \right\} \\ = \exp \left[ \mathbb {E} \log \left\{\left(\det  L\right) \det  \left(L ^ {\top}\right) \right\} \right], \\ \end{array}
$$

$$
\begin{array}{l} = \exp \left\{\mathbb {E} \log \det  (L L ^ {\top}) \right\} \\ = \exp \left\{\mathbb {E} \log \det  (P) \right\}, \\ \end{array}
$$

[37] S. SRA, Positive definite matrices and the  $S$ -divergence, Proc. Amer. Math. Soc., 144 (2016), pp. 2787-2797.

[38] B. VANDEREYCKEN, P.-A. ABSIL, AND S. VANDEWALLE, A Riemannian geometry with complete geodesics for the set of positive semidefinite matrices of fixed rank, IMA J. Numer. Anal., 33 (2013), pp. 481-514.

[39] Z. WANG, B. C. VEMURI, Y. CHEN, AND T. H. MARECI, A constrained variational principle for direct estimation and smoothing of the diffusion tensor field from complex DWI, IEEE Trans. Medical Imaging, 23 (2004), pp. 930-939.


\begin{proposition}[Proposition 10]
PROPOSITION 10. Suppose  $L$  is a random element on  $\mathcal{L}_{+}$  and  $P$  is a random element on  $S_{m}^{+}$ . Suppose for some fixed  $A \in \mathcal{L}_{+}$  and  $B \in S_{m}^{+}$  such that  $\mathbb{E}d_{\mathcal{L}_{+}}^{2}(A,L) < \infty$  and  $\mathbb{E}d_{S_m^+}^2 (B,P) < \infty$ . Then the Fréchet mean of  $L$  is given by

$$
\mathbb {E} L = \mathbb {E} \lfloor L \rfloor + \exp \left\{\mathbb {E} \log \mathbb {D} (L) \right\}, \tag {4.1}
$$
\end{proposition}

\begin{proof}
Proof. Let  $F(R) = \mathbb{E}d_{\mathcal{L}_+}^2 (R,L)$ . Then

$$
F (R) = \mathbb {E} \| \lfloor R \rfloor - \lfloor L \rfloor \| _ {\mathrm {F}} ^ {2} + \mathbb {E} \| \log \mathbb {D} (R) - \log \mathbb {D} (L) \| _ {\mathrm {F}} ^ {2} := F _ {1} (\lfloor R \rfloor) + F _ {2} (\mathbb {D} (R)).
$$

To minimize  $F$ , we can separately minimize  $F_{1}$  and  $F_{2}$ , due to two reasons: (1)  $F_{1}$  involves  $\lfloor R\rfloor$  only while  $F_{2}$  involves  $\mathbb{D}(R)$ , and (2)  $\lfloor R\rfloor$  and  $\mathbb{D}(R)$  are disjoint and can vary independently. For  $F_{1}$ , we first note that the condition  $\mathbb{E}d_{\mathcal{L}_{+}}^{2}(A,L) < \infty$  ensures that  $\mathbb{E}\| \lfloor L\rfloor \|_{\mathrm{F}}^2 < \infty$ , and hence the mean  $\mathbb{E}\lfloor L\rfloor$  is well defined. Also, it can be checked that this mean minimizes  $F_{1}$ .

For  $F_{2}$ , we note that

$$
\begin{array}{l} F _ {2} (\mathbb {D} (R)) = \sum_ {j = 1} ^ {m} \mathbb {E} (\log \mathbb {D} (R) - \log \mathbb {D} (L)) _ {j j} ^ {2} \\ = \sum_ {j = 1} ^ {m} \mathbb {E} \left\{\log R _ {j j} - \log L _ {j j} \right\} ^ {2} := \sum_ {j = 1} ^ {m} F _ {2 j} \left(R _ {j j}\right). \\ \end{array}
$$

Again, these  $m$  components  $F_{2j}$  can be optimized separately. From the condition  $\mathbb{E}d_{\mathcal{L}_+}^2 (A,L) < \infty$  we deduce that  $\mathbb{E}(\log L_{jj})^2 < \infty$ . Thus,  $\mathbb{E}\log L_{jj}$  is well defined and minimizes  $F_{2j}(e^{x})$  with respect to  $x$ . Therefore,  $\exp \{\mathbb{E}\log L_{jj}\}$  minimizes  $F_{2j}(x)$  for each  $j = 1,\ldots,m$ . In matrix form, this is equivalent to that  $\exp \{\mathbb{E}\log \mathbb{D}(L)\}$  minimizes  $F_{2}(D)$  when  $D$  is constrained to be diagonal. Thus, combined with the optimizer for  $F_{1}$ , it establishes (4.1). Then (4.2) follows from the fact that  $\mathcal{S}(L) = LL^{\top}$  is an isometry.

Finally, we establish the following useful relation between the determinant of the Log-Cholesky mean and the mean of the logarithmic determinant of a random SPD matrix.

\end{proof}

\paragraph{Background.}
-  $\operatorname{det}(X) = \prod_{j=1}^{m} X_{jj}$  for  $X \in \mathcal{L}$ .

These properties show that both  $\mathcal{L}$  and  $\mathcal{L}_{+}$  are closed under matrix addition and multiplication and that the operator  $\mathbb{D}$  interacts well with these operations.

Recall that  $S_{m}^{+}$  is defined as the collection of  $m \times m$  SPD matrices. We denote the space of  $m \times m$  symmetric space by  $S_{m}$ . Symmetric matrices and SPD matrices possess numerous algebraic and analytic properties that are well documented in [6]. Below are some of them to be used in what follows.

- For any invertible matrix  $X$ , the matrix  $XX^{\top}$  is an SPD matrix.

-  $\exp(S)$  is an SPD matrix for a symmetric matrix  $S$ , while  $\log(P)$  is a symmetric metric for an SPD matrix  $P$ .

- Diagonal elements of an SPD matrix are all positive. This can be seen from the fact that  $P_{jj} = e_j^\top Pe_j > 0$  for  $P \in S_m^+$ , where  $e_j$  is the unit vector with 1 at the  $j$ th coordinate and 0 elsewhere.

-  $\exp(S)$  is an SPD matrix for a symmetric matrix  $S$ , while  $\log(P)$  is a symmetric metric for an SPD matrix  $P$ .

- Diagonal elements of an SPD matrix are all positive. This can be seen from the fact that  $P_{jj} = e_j^\top Pe_j > 0$  for  $P \in S_m^+$ , where  $e_j$  is the unit vector with 1 at the  $j$ th coordinate and 0 elsewhere.

2.2. Cholesky decomposition. Cholesky decomposition, named after André-Louis Cholesky, represents a real  $m \times m$  SPD matrix  $P$  as a product of a lower triangular matrix  $L$  and its transpose, i.e.,  $P = LL^{\top}$ . If the diagonal elements of  $L$  are restricted to be positive, then the decomposition is unique according to Theorem 4.2.5 of [16]. Such lower triangular matrix, denoted by  $\mathcal{L}(P)$ , is called the Cholesky factor of  $P$ . Since in addition  $L = \mathcal{L}(LL^{\top})$  for each  $L \in \mathcal{L}_{+}$ , the map  $\mathcal{L}: S_{m}^{+} \to \mathcal{L}_{+}$  is bijective. In other words, there is one-to-one correspondence between SPD matrices and lower triangular matrices whose diagonal elements are all positive.

3.1. Riemannian geometry on Cholesky spaces. For matrices in  $\mathcal{L}_{+}$ , as off-diagonal elements in the lower triangular part are unconstrained while diagonal ones are restricted to be positive,  $\mathcal{L}_{+}$  can be parameterized by  $\mathcal{L} \ni X \coloneqq \varphi(X) \in \mathcal{L}_{+}$  in the way that  $(\varphi(X))_{ij} = X_{ij}$  if  $i \neq j$  and  $(\varphi(X))_{jj} = \exp(X_{jj})$ . This motivates us to, respectively, endow the unconstrained part  $\lfloor \cdot \rfloor$  and the positive diagonal part  $\mathbb{D}(\cdot)$  with a different metric and then combine them into a Riemannian metric on  $\mathcal{L}_{+}$  as follows. First, we note that the tangent space of  $\mathcal{L}_{+}$  at a given  $L \in \mathcal{L}_{+}$  is identified with the linear space  $\mathcal{L}$ . For such tangent space, we treat the strict lower triangular space  $\lfloor \mathcal{L} \rfloor \coloneqq \{\lfloor X \rfloor : X \in \mathcal{L}\}$  as the Euclidean space  $\mathbb{R}^{m(m-1)/2}$  with the usual Frobenius inner product  $\langle X, Y \rangle_{\mathrm{F}} = \sum_{i,j=1}^{m} X_{ij} Y_{ij}$  for all  $X, Y \in \lfloor \mathcal{L} \rfloor$ . For the diagonal part  $\mathbb{D}(\mathcal{L}) \coloneqq \{\mathbb{D}(Z) : Z \in \mathcal{L}\}$ , we equipped it with a different inner product defined by  $\langle \mathbb{D}(L)^{-1} \mathbb{D}(X), \mathbb{D}(L)^{-1} \mathbb{D}(Y) \rangle_{\mathrm{F}}$ . Finally, combining these two components together, we define a metric  $\tilde{g}$  for tangent spaces  $T_L \mathcal{L}_{+}$  (identified with  $\mathcal{L}$ ) by

$$
\begin{array}{l} \tilde {g} _ {L} (X, Y) = \langle \lfloor X \rfloor , \lfloor Y \rfloor \rangle_ {\mathrm {F}} + \langle \mathbb {D} (L) ^ {- 1} \mathbb {D} (X), \mathbb {D} (L) ^ {- 1} \mathbb {D} (Y) \rangle_ {\mathrm {F}} \\ = \sum_ {i > j} X _ {i j} Y _ {i j} + \sum_ {j = 1} ^ {m} X _ {j j} Y _ {j j} L _ {j j} ^ {- 2}. \\ \end{array}
$$

It is straightforward to show that the space  $\mathcal{L}_{+}$ , equipped with the metric  $\tilde{g}$ , is a Riemannian manifold. We begin with geodesics on the manifold  $(\mathcal{L}_{+},\tilde{g})$  to investigate its basic properties.

$$
L ^ {- 1} W L ^ {- \top} = X ^ {\top} L ^ {- \top} + L ^ {- 1} X = L ^ {- 1} X + (L ^ {- 1} X) ^ {\top}.
$$

Note that  $L^{-1}X$  is also a lower triangular matrix, and the matrix on the left-hand side is symmetric, we deduce that  $(L^{-1}WL^{-\top})_{\frac{1}{2}} = L^{-1}X$ , which gives rise to  $X = L(L^{-1}WL^{-\top})_{\frac{1}{2}}$ . The linear map  $(D_L\mathcal{S})(X) = LX^\top +XL^\top$  is one-to-one, since from the above derivation,  $(D_L\mathcal{S})(X) = 0$  if and only if  $X = 0$ .

Given the above proposition, the manifold map  $\mathcal{S}$ , which is exactly the inverse of the Cholesky map  $\mathcal{L}$  discussed in subsection 2.2, induces a Riemannian metric on  $S_{m}^{+}$ , denoted by  $g$  and called Log-Cholesky metric, given by

Note that  $L^{-1}X$  is also a lower triangular matrix, and the matrix on the left-hand side is symmetric, we deduce that  $(L^{-1}WL^{-\top})_{\frac{1}{2}} = L^{-1}X$ , which gives rise to  $X = L(L^{-1}WL^{-\top})_{\frac{1}{2}}$ . The linear map  $(D_L\mathcal{S})(X) = LX^\top +XL^\top$  is one-to-one, since from the above derivation,  $(D_L\mathcal{S})(X) = 0$  if and only if  $X = 0$ .

Given the above proposition, the manifold map  $\mathcal{S}$ , which is exactly the inverse of the Cholesky map  $\mathcal{L}$  discussed in subsection 2.2, induces a Riemannian metric on  $S_{m}^{+}$ , denoted by  $g$  and called Log-Cholesky metric, given by

$$
g _ {P} (W, V) = \tilde {g} _ {L} \left(L \left(L ^ {- 1} W L ^ {- \top}\right) _ {\frac {1}{2}}, L \left(L ^ {- 1} V L ^ {- \top}\right) _ {\frac {1}{2}}\right), \tag {3.1}
$$

$$
\gamma_ {P, W} (t) = \mathcal {S} (\tilde {\gamma} _ {L, X} (t)) = \tilde {\gamma} _ {L, X} (t) \tilde {\gamma} _ {L, X} ^ {\top} (t),
$$

where  $X = L(L^{-1}WL^{-\top})_{\frac{1}{2}}$  and  $W \in T_P\mathcal{S}_m^+$ . Similarly, the Riemannian exponential at  $P$  is given by

$$
\mathrm {E x p} _ {P} W = \mathcal {S} (\widetilde {\mathrm {E x p}} _ {L} X) = (\widetilde {\mathrm {E x p}} _ {L} X) (\widetilde {\mathrm {E x p}} _ {L} X) ^ {\top},
$$

where  $X = L(L^{-1}WL^{-\top})_{\frac{1}{2}}$  and  $W \in T_P\mathcal{S}_m^+$ . Similarly, the Riemannian exponential at  $P$  is given by

$$
\mathrm {E x p} _ {P} W = \mathcal {S} (\widetilde {\mathrm {E x p}} _ {L} X) = (\widetilde {\mathrm {E x p}} _ {L} X) (\widetilde {\mathrm {E x p}} _ {L} X) ^ {\top},
$$

while the geodesic between  $P$  and  $Q$  is characterized by

3.3. Lie group structure and bi-invariant metrics. We define an operator  $\odot$  on  $\mathcal{L}$  by

$$
X \circledast Y = \left\lfloor X \right\rfloor + \left\lfloor Y \right\rfloor + \mathbb {D} (X) \mathbb {D} (Y).
$$

Note that  $\mathcal{L}_{+} \subset \mathcal{L}$ . Moreover, if  $L, K \in \mathcal{L}_{+}$ , then  $L \circledast K \in \mathcal{L}_{+}$ . It is not difficult to see that  $\circledast$  is a smooth commutative group operation on the manifold  $\mathcal{L}_{+}$ , where the

$$
X \circledast Y = \left\lfloor X \right\rfloor + \left\lfloor Y \right\rfloor + \mathbb {D} (X) \mathbb {D} (Y).
$$

Note that  $\mathcal{L}_{+} \subset \mathcal{L}$ . Moreover, if  $L, K \in \mathcal{L}_{+}$ , then  $L \circledast K \in \mathcal{L}_{+}$ . It is not difficult to see that  $\circledast$  is a smooth commutative group operation on the manifold  $\mathcal{L}_{+}$ , where the

# LOG-CHOLESKY METRIC

inverse of  $L$ , denoted by  $L_{\odot}^{-1}$ , is  $\mathbb{D}(L)^{-1} - \lfloor L\rfloor$ . The left translation by  $A\in \mathcal{L}_{+}$  is denoted by  $\ell_A:B\mapsto A\odot B$ . One can check that the differential of this operation at  $L\in \mathcal{L}_{+}$  is

$$
D _ {L} \ell_ {A}: X \mapsto \lfloor X \rfloor + \mathbb {D} (A) \mathbb {D} (X), \tag {3.2}
$$

where it is noted that the differential  $D_{L}\ell_{A}$  does not depend on  $L$ . Given the above expression, one can find that

$$
D _ {L} \ell_ {A}: X \mapsto \lfloor X \rfloor + \mathbb {D} (A) \mathbb {D} (X), \tag {3.2}
$$

where it is noted that the differential  $D_{L}\ell_{A}$  does not depend on  $L$ . Given the above expression, one can find that

$$
\begin{array}{l} \tilde {g} _ {A \odot L} ((D _ {L} \ell_ {A}) (X), (D _ {L} \ell_ {A}) (Y)) \\ = \tilde {g} _ {A \odot L} (\left\lfloor X \right\rfloor + \mathbb {D} (A) \mathbb {D} (X), \left\lfloor Y \right\rfloor + \mathbb {D} (A) \mathbb {D} (Y)) \\ = \tilde {g} _ {L} (X, Y). \\ \end{array}
$$

The group operator  $\odot$  and maps  $\mathcal{S}$  and  $\mathcal{L}$  together induce a smooth operation  $\otimes$  on  $S_{m}^{+}$ , defined by

$$
\begin{array}{l} P \circledast Q = \mathcal {S} (\mathcal {L} (P) \circledast \mathcal {L} (Q)) \\ = \left(\mathcal {L} (P) \circledast \mathcal {L} (Q)\right) \left(\mathcal {L} (P) \circledast \mathcal {L} (Q)\right) ^ {\top} \quad \text {f o r} P, Q \in S _ {m} ^ {+}. \\ \end{array}
$$

In addition, both  $\mathcal{L}$  and  $\mathcal{S}$  are Riemannian isometries and group isomorphisms between Lie groups  $(\mathcal{L}_{+},\tilde{g},\circledast)$  and  $(\mathcal{S}_m^+,g,\bullet)$ .

$$
\begin{array}{l} \det  \mathbb {E} L = \exp [ \operatorname {t r} \{\log \mathbb {D} (\mathbb {E} L) \} ] = \exp [ \operatorname {t r} \{\mathbb {E} \log \mathbb {D} (L) \} ] \\ = \exp \left[ \mathbb {E} \operatorname {t r} \left\{\log \mathbb {D} (L) \right\} \right] = \exp \left\{\mathbb {E} \log (\det  L) \right\}, \\ \end{array}
$$

which proves the case of  $(\mathcal{L}_{+},\tilde{g})$

For the case that  $Q$  is a random element on  $(\mathcal{S}_m^+, g)$ , we denote it by  $P$  and observe that  $\operatorname{det}(AB) = (\operatorname{det} A)(\operatorname{det} B)$  for any square matrices  $A$  and  $B$ . Let  $L = \mathcal{L}(P)$ . Then, one has

which proves the case of  $(\mathcal{L}_{+},\tilde{g})$

For the case that  $Q$  is a random element on  $(\mathcal{S}_m^+, g)$ , we denote it by  $P$  and observe that  $\operatorname{det}(AB) = (\operatorname{det} A)(\operatorname{det} B)$  for any square matrices  $A$  and  $B$ . Let  $L = \mathcal{L}(P)$ . Then, one has

$$
\begin{array}{l} \det  \mathbb {E} P = \det  \left\{\left(\mathbb {E} L\right) \left(\mathbb {E} L\right) ^ {\top} \right\} \\ = \left\{\det  (\mathbb {E} L) \right\} \left\{\det  (\mathbb {E} L) ^ {\top} \right\} \\ = \exp \left\{\mathbb {E} \log (\det  L) \right\} \exp \left\{\mathbb {E} \log (\det  L ^ {\top}) \right\} \\ = \exp \left\{\mathbb {E} \left(\log \det  L + \log \det  L ^ {\top}\right) \right\} \\ = \exp \left[ \mathbb {E} \log \left\{\left(\det  L\right) \det  \left(L ^ {\top}\right) \right\} \right], \\ \end{array}
$$

For the case that  $Q$  is a random element on  $(\mathcal{S}_m^+, g)$ , we denote it by  $P$  and observe that  $\operatorname{det}(AB) = (\operatorname{det} A)(\operatorname{det} B)$  for any square matrices  $A$  and  $B$ . Let  $L = \mathcal{L}(P)$ . Then, one has

$$
\begin{array}{l} \det  \mathbb {E} P = \det  \left\{\left(\mathbb {E} L\right) \left(\mathbb {E} L\right) ^ {\top} \right\} \\ = \left\{\det  (\mathbb {E} L) \right\} \left\{\det  (\mathbb {E} L) ^ {\top} \right\} \\ = \exp \left\{\mathbb {E} \log (\det  L) \right\} \exp \left\{\mathbb {E} \log (\det  L ^ {\top}) \right\} \\ = \exp \left\{\mathbb {E} \left(\log \det  L + \log \det  L ^ {\top}\right) \right\} \\ = \exp \left[ \mathbb {E} \log \left\{\left(\det  L\right) \det  \left(L ^ {\top}\right) \right\} \right], \\ \end{array}
$$

$$
\begin{array}{l} = \exp \left\{\mathbb {E} \log \det  (L L ^ {\top}) \right\} \\ = \exp \left\{\mathbb {E} \log \det  (P) \right\}, \\ \end{array}
$$


\begin{proposition}[Proposition 11]
PROPOSITION 11. If the Fréchet mean of a random element  $Q$  on  $(\mathcal{L}_{+},\tilde{g})$  or  $(S_{m}^{+},g)$  exists, then  $\log \det(\mathbb{E}Q) = \mathbb{E}\log(\det Q)$ .
\end{proposition}

\begin{proof}
Proof. For the case that  $Q$  is a random element on  $(\mathcal{L}_{+},\tilde{g})$ , we denote it by  $L$  and observe that in (4.1)  $\mathbb{E}L$  is a lower triangular matrix and  $\exp \left\{\mathbb{E}\log \mathbb{D}(L)\right\}$  is its diagonal part. For a triangular matrix, its determinant is the product of diagonal elements. Thus,

$$
\det L = \prod_ {j = 1} ^ {m} L _ {j j} = \exp [ \mathrm {t r} \{\log \mathbb {D} (L) \} ],
$$

or equivalently,  $\log \det L = \operatorname{tr}\{\log \mathbb{D}(L)\}$ . The above observations imply that

$$
\begin{array}{l} \det  \mathbb {E} L = \exp [ \operatorname {t r} \{\log \mathbb {D} (\mathbb {E} L) \} ] = \exp [ \operatorname {t r} \{\mathbb {E} \log \mathbb {D} (L) \} ] \\ = \exp \left[ \mathbb {E} \operatorname {t r} \left\{\log \mathbb {D} (L) \right\} \right] = \exp \left\{\mathbb {E} \log (\det  L) \right\}, \\ \end{array}
$$

which proves the case of  $(\mathcal{L}_{+},\tilde{g})$

For the case that  $Q$  is a random element on  $(\mathcal{S}_m^+, g)$ , we denote it by  $P$  and observe that  $\operatorname{det}(AB) = (\operatorname{det} A)(\operatorname{det} B)$  for any square matrices  $A$  and  $B$ . Let  $L = \mathcal{L}(P)$ . Then, one has

$$
\begin{array}{l} \det  \mathbb {E} P = \det  \left\{\left(\mathbb {E} L\right) \left(\mathbb {E} L\right) ^ {\top} \right\} \\ = \left\{\det  (\mathbb {E} L) \right\} \left\{\det  (\mathbb {E} L) ^ {\top} \right\} \\ = \exp \left\{\mathbb {E} \log (\det  L) \right\} \exp \left\{\mathbb {E} \log (\det  L ^ {\top}) \right\} \\ = \exp \left\{\mathbb {E} \left(\log \det  L + \log \det  L ^ {\top}\right) \right\} \\ = \exp \left[ \mathbb {E} \log \left\{\left(\det  L\right) \det  \left(L ^ {\top}\right) \right\} \right], \\ \end{array}
$$

$$
\begin{array}{l} = \exp \left\{\mathbb {E} \log \det  (L L ^ {\top}) \right\} \\ = \exp \left\{\mathbb {E} \log \det  (P) \right\}, \\ \end{array}
$$

which establishes the statement for the case of  $(S_m^+, g)$ .

\end{proof}

\paragraph{Background.}
-  $\operatorname{det}(X) = \prod_{j=1}^{m} X_{jj}$  for  $X \in \mathcal{L}$ .

These properties show that both  $\mathcal{L}$  and  $\mathcal{L}_{+}$  are closed under matrix addition and multiplication and that the operator  $\mathbb{D}$  interacts well with these operations.

Recall that  $S_{m}^{+}$  is defined as the collection of  $m \times m$  SPD matrices. We denote the space of  $m \times m$  symmetric space by  $S_{m}$ . Symmetric matrices and SPD matrices possess numerous algebraic and analytic properties that are well documented in [6]. Below are some of them to be used in what follows.

-  $\exp(S)$  is an SPD matrix for a symmetric matrix  $S$ , while  $\log(P)$  is a symmetric metric for an SPD matrix  $P$ .

- Diagonal elements of an SPD matrix are all positive. This can be seen from the fact that  $P_{jj} = e_j^\top Pe_j > 0$  for  $P \in S_m^+$ , where  $e_j$  is the unit vector with 1 at the  $j$ th coordinate and 0 elsewhere.

2.2. Cholesky decomposition. Cholesky decomposition, named after André-Louis Cholesky, represents a real  $m \times m$  SPD matrix  $P$  as a product of a lower triangular matrix  $L$  and its transpose, i.e.,  $P = LL^{\top}$ . If the diagonal elements of  $L$  are restricted to be positive, then the decomposition is unique according to Theorem 4.2.5 of [16]. Such lower triangular matrix, denoted by  $\mathcal{L}(P)$ , is called the Cholesky factor of  $P$ . Since in addition  $L = \mathcal{L}(LL^{\top})$  for each  $L \in \mathcal{L}_{+}$ , the map  $\mathcal{L}: S_{m}^{+} \to \mathcal{L}_{+}$  is bijective. In other words, there is one-to-one correspondence between SPD matrices and lower triangular matrices whose diagonal elements are all positive.

3.1. Riemannian geometry on Cholesky spaces. For matrices in  $\mathcal{L}_{+}$ , as off-diagonal elements in the lower triangular part are unconstrained while diagonal ones are restricted to be positive,  $\mathcal{L}_{+}$  can be parameterized by  $\mathcal{L} \ni X \coloneqq \varphi(X) \in \mathcal{L}_{+}$  in the way that  $(\varphi(X))_{ij} = X_{ij}$  if  $i \neq j$  and  $(\varphi(X))_{jj} = \exp(X_{jj})$ . This motivates us to, respectively, endow the unconstrained part  $\lfloor \cdot \rfloor$  and the positive diagonal part  $\mathbb{D}(\cdot)$  with a different metric and then combine them into a Riemannian metric on  $\mathcal{L}_{+}$  as follows. First, we note that the tangent space of  $\mathcal{L}_{+}$  at a given  $L \in \mathcal{L}_{+}$  is identified with the linear space  $\mathcal{L}$ . For such tangent space, we treat the strict lower triangular space  $\lfloor \mathcal{L} \rfloor \coloneqq \{\lfloor X \rfloor : X \in \mathcal{L}\}$  as the Euclidean space  $\mathbb{R}^{m(m-1)/2}$  with the usual Frobenius inner product  $\langle X, Y \rangle_{\mathrm{F}} = \sum_{i,j=1}^{m} X_{ij} Y_{ij}$  for all  $X, Y \in \lfloor \mathcal{L} \rfloor$ . For the diagonal part  $\mathbb{D}(\mathcal{L}) \coloneqq \{\mathbb{D}(Z) : Z \in \mathcal{L}\}$ , we equipped it with a different inner product defined by  $\langle \mathbb{D}(L)^{-1} \mathbb{D}(X), \mathbb{D}(L)^{-1} \mathbb{D}(Y) \rangle_{\mathrm{F}}$ . Finally, combining these two components together, we define a metric  $\tilde{g}$  for tangent spaces  $T_L \mathcal{L}_{+}$  (identified with  $\mathcal{L}$ ) by

$$
\begin{array}{l} \tilde {g} _ {L} (X, Y) = \langle \lfloor X \rfloor , \lfloor Y \rfloor \rangle_ {\mathrm {F}} + \langle \mathbb {D} (L) ^ {- 1} \mathbb {D} (X), \mathbb {D} (L) ^ {- 1} \mathbb {D} (Y) \rangle_ {\mathrm {F}} \\ = \sum_ {i > j} X _ {i j} Y _ {i j} + \sum_ {j = 1} ^ {m} X _ {j j} Y _ {j j} L _ {j j} ^ {- 2}. \\ \end{array}
$$

It is straightforward to show that the space  $\mathcal{L}_{+}$ , equipped with the metric  $\tilde{g}$ , is a Riemannian manifold. We begin with geodesics on the manifold  $(\mathcal{L}_{+},\tilde{g})$  to investigate its basic properties.

$$
L ^ {- 1} W L ^ {- \top} = X ^ {\top} L ^ {- \top} + L ^ {- 1} X = L ^ {- 1} X + (L ^ {- 1} X) ^ {\top}.
$$

Note that  $L^{-1}X$  is also a lower triangular matrix, and the matrix on the left-hand side is symmetric, we deduce that  $(L^{-1}WL^{-\top})_{\frac{1}{2}} = L^{-1}X$ , which gives rise to  $X = L(L^{-1}WL^{-\top})_{\frac{1}{2}}$ . The linear map  $(D_L\mathcal{S})(X) = LX^\top +XL^\top$  is one-to-one, since from the above derivation,  $(D_L\mathcal{S})(X) = 0$  if and only if  $X = 0$ .

Given the above proposition, the manifold map  $\mathcal{S}$ , which is exactly the inverse of the Cholesky map  $\mathcal{L}$  discussed in subsection 2.2, induces a Riemannian metric on  $S_{m}^{+}$ , denoted by  $g$  and called Log-Cholesky metric, given by

Note that  $L^{-1}X$  is also a lower triangular matrix, and the matrix on the left-hand side is symmetric, we deduce that  $(L^{-1}WL^{-\top})_{\frac{1}{2}} = L^{-1}X$ , which gives rise to  $X = L(L^{-1}WL^{-\top})_{\frac{1}{2}}$ . The linear map  $(D_L\mathcal{S})(X) = LX^\top +XL^\top$  is one-to-one, since from the above derivation,  $(D_L\mathcal{S})(X) = 0$  if and only if  $X = 0$ .

Given the above proposition, the manifold map  $\mathcal{S}$ , which is exactly the inverse of the Cholesky map  $\mathcal{L}$  discussed in subsection 2.2, induces a Riemannian metric on  $S_{m}^{+}$ , denoted by  $g$  and called Log-Cholesky metric, given by

$$
g _ {P} (W, V) = \tilde {g} _ {L} \left(L \left(L ^ {- 1} W L ^ {- \top}\right) _ {\frac {1}{2}}, L \left(L ^ {- 1} V L ^ {- \top}\right) _ {\frac {1}{2}}\right), \tag {3.1}
$$

where  $X = L(L^{-1}WL^{-\top})_{\frac{1}{2}}$  and  $W \in T_P\mathcal{S}_m^+$ . Similarly, the Riemannian exponential at  $P$  is given by

$$
\mathrm {E x p} _ {P} W = \mathcal {S} (\widetilde {\mathrm {E x p}} _ {L} X) = (\widetilde {\mathrm {E x p}} _ {L} X) (\widetilde {\mathrm {E x p}} _ {L} X) ^ {\top},
$$

while the geodesic between  $P$  and  $Q$  is characterized by

inverse of  $L$ , denoted by  $L_{\odot}^{-1}$ , is  $\mathbb{D}(L)^{-1} - \lfloor L\rfloor$ . The left translation by  $A\in \mathcal{L}_{+}$  is denoted by  $\ell_A:B\mapsto A\odot B$ . One can check that the differential of this operation at  $L\in \mathcal{L}_{+}$  is

$$
D _ {L} \ell_ {A}: X \mapsto \lfloor X \rfloor + \mathbb {D} (A) \mathbb {D} (X), \tag {3.2}
$$

where it is noted that the differential  $D_{L}\ell_{A}$  does not depend on  $L$ . Given the above expression, one can find that

$$
D _ {L} \ell_ {A}: X \mapsto \lfloor X \rfloor + \mathbb {D} (A) \mathbb {D} (X), \tag {3.2}
$$

where it is noted that the differential  $D_{L}\ell_{A}$  does not depend on  $L$ . Given the above expression, one can find that

$$
\begin{array}{l} \tilde {g} _ {A \odot L} ((D _ {L} \ell_ {A}) (X), (D _ {L} \ell_ {A}) (Y)) \\ = \tilde {g} _ {A \odot L} (\left\lfloor X \right\rfloor + \mathbb {D} (A) \mathbb {D} (X), \left\lfloor Y \right\rfloor + \mathbb {D} (A) \mathbb {D} (Y)) \\ = \tilde {g} _ {L} (X, Y). \\ \end{array}
$$

$$
\begin{array}{l} P \circledast Q = \mathcal {S} (\mathcal {L} (P) \circledast \mathcal {L} (Q)) \\ = \left(\mathcal {L} (P) \circledast \mathcal {L} (Q)\right) \left(\mathcal {L} (P) \circledast \mathcal {L} (Q)\right) ^ {\top} \quad \text {f o r} P, Q \in S _ {m} ^ {+}. \\ \end{array}
$$

In addition, both  $\mathcal{L}$  and  $\mathcal{S}$  are Riemannian isometries and group isomorphisms between Lie groups  $(\mathcal{L}_{+},\tilde{g},\circledast)$  and  $(\mathcal{S}_m^+,g,\bullet)$ .

THEOREM 5. The space  $(S_m^+, \otimes)$  is an abelian Lie group. Moreover, the metric  $g$  defined in (3.1) is a bi-invariant metric on  $(S_m^+, \otimes)$ .

$$
\begin{array}{l} \det  \mathbb {E} L = \exp [ \operatorname {t r} \{\log \mathbb {D} (\mathbb {E} L) \} ] = \exp [ \operatorname {t r} \{\mathbb {E} \log \mathbb {D} (L) \} ] \\ = \exp \left[ \mathbb {E} \operatorname {t r} \left\{\log \mathbb {D} (L) \right\} \right] = \exp \left\{\mathbb {E} \log (\det  L) \right\}, \\ \end{array}
$$

which proves the case of  $(\mathcal{L}_{+},\tilde{g})$

For the case that  $Q$  is a random element on  $(\mathcal{S}_m^+, g)$ , we denote it by  $P$  and observe that  $\operatorname{det}(AB) = (\operatorname{det} A)(\operatorname{det} B)$  for any square matrices  $A$  and  $B$ . Let  $L = \mathcal{L}(P)$ . Then, one has

which proves the case of  $(\mathcal{L}_{+},\tilde{g})$

For the case that  $Q$  is a random element on  $(\mathcal{S}_m^+, g)$ , we denote it by  $P$  and observe that  $\operatorname{det}(AB) = (\operatorname{det} A)(\operatorname{det} B)$  for any square matrices  $A$  and  $B$ . Let  $L = \mathcal{L}(P)$ . Then, one has

$$
\begin{array}{l} \det  \mathbb {E} P = \det  \left\{\left(\mathbb {E} L\right) \left(\mathbb {E} L\right) ^ {\top} \right\} \\ = \left\{\det  (\mathbb {E} L) \right\} \left\{\det  (\mathbb {E} L) ^ {\top} \right\} \\ = \exp \left\{\mathbb {E} \log (\det  L) \right\} \exp \left\{\mathbb {E} \log (\det  L ^ {\top}) \right\} \\ = \exp \left\{\mathbb {E} \left(\log \det  L + \log \det  L ^ {\top}\right) \right\} \\ = \exp \left[ \mathbb {E} \log \left\{\left(\det  L\right) \det  \left(L ^ {\top}\right) \right\} \right], \\ \end{array}
$$

For the case that  $Q$  is a random element on  $(\mathcal{S}_m^+, g)$ , we denote it by  $P$  and observe that  $\operatorname{det}(AB) = (\operatorname{det} A)(\operatorname{det} B)$  for any square matrices  $A$  and  $B$ . Let  $L = \mathcal{L}(P)$ . Then, one has

$$
\begin{array}{l} \det  \mathbb {E} P = \det  \left\{\left(\mathbb {E} L\right) \left(\mathbb {E} L\right) ^ {\top} \right\} \\ = \left\{\det  (\mathbb {E} L) \right\} \left\{\det  (\mathbb {E} L) ^ {\top} \right\} \\ = \exp \left\{\mathbb {E} \log (\det  L) \right\} \exp \left\{\mathbb {E} \log (\det  L ^ {\top}) \right\} \\ = \exp \left\{\mathbb {E} \left(\log \det  L + \log \det  L ^ {\top}\right) \right\} \\ = \exp \left[ \mathbb {E} \log \left\{\left(\det  L\right) \det  \left(L ^ {\top}\right) \right\} \right], \\ \end{array}
$$

$$
\begin{array}{l} = \exp \left\{\mathbb {E} \log \det  (L L ^ {\top}) \right\} \\ = \exp \left\{\mathbb {E} \log \det  (P) \right\}, \\ \end{array}
$$


\section*{Corollarys}

\begin{corollary}[Corollary 12]
COROLLARY 12. For  $L_{1},\ldots ,L_{n}\in \mathcal{L}_{+}$  , one has

$$
\mathbb {E} _ {n} \left(L _ {1}, \dots , L _ {n}\right) = \frac {1}{n} \sum_ {i = 1} ^ {n} \left\lfloor L _ {i} \right\rfloor + \exp \left\{n ^ {- 1} \sum_ {i = 1} ^ {n} \log \mathbb {D} \left(L _ {i}\right) \right\}, \tag {4.3}
$$
\end{corollary}

\paragraph{Background.}
-  $\exp(S)$  is an SPD matrix for a symmetric matrix  $S$ , while  $\log(P)$  is a symmetric metric for an SPD matrix  $P$ .

- Diagonal elements of an SPD matrix are all positive. This can be seen from the fact that  $P_{jj} = e_j^\top Pe_j > 0$  for  $P \in S_m^+$ , where  $e_j$  is the unit vector with 1 at the  $j$ th coordinate and 0 elsewhere.

2.2. Cholesky decomposition. Cholesky decomposition, named after André-Louis Cholesky, represents a real  $m \times m$  SPD matrix  $P$  as a product of a lower triangular matrix  $L$  and its transpose, i.e.,  $P = LL^{\top}$ . If the diagonal elements of  $L$  are restricted to be positive, then the decomposition is unique according to Theorem 4.2.5 of [16]. Such lower triangular matrix, denoted by  $\mathcal{L}(P)$ , is called the Cholesky factor of  $P$ . Since in addition  $L = \mathcal{L}(LL^{\top})$  for each  $L \in \mathcal{L}_{+}$ , the map  $\mathcal{L}: S_{m}^{+} \to \mathcal{L}_{+}$  is bijective. In other words, there is one-to-one correspondence between SPD matrices and lower triangular matrices whose diagonal elements are all positive.

inverse of  $L$ , denoted by  $L_{\odot}^{-1}$ , is  $\mathbb{D}(L)^{-1} - \lfloor L\rfloor$ . The left translation by  $A\in \mathcal{L}_{+}$  is denoted by  $\ell_A:B\mapsto A\odot B$ . One can check that the differential of this operation at  $L\in \mathcal{L}_{+}$  is

$$
D _ {L} \ell_ {A}: X \mapsto \lfloor X \rfloor + \mathbb {D} (A) \mathbb {D} (X), \tag {3.2}
$$

where it is noted that the differential  $D_{L}\ell_{A}$  does not depend on  $L$ . Given the above expression, one can find that

$$
D _ {L} \ell_ {A}: X \mapsto \lfloor X \rfloor + \mathbb {D} (A) \mathbb {D} (X), \tag {3.2}
$$

where it is noted that the differential  $D_{L}\ell_{A}$  does not depend on  $L$ . Given the above expression, one can find that

$$
\begin{array}{l} \tilde {g} _ {A \odot L} ((D _ {L} \ell_ {A}) (X), (D _ {L} \ell_ {A}) (Y)) \\ = \tilde {g} _ {A \odot L} (\left\lfloor X \right\rfloor + \mathbb {D} (A) \mathbb {D} (X), \left\lfloor Y \right\rfloor + \mathbb {D} (A) \mathbb {D} (Y)) \\ = \tilde {g} _ {L} (X, Y). \\ \end{array}
$$


\begin{corollary}[Corollary 13]
COROLLARY 13. For  $L_{1},\ldots ,L_{n}\in \mathcal{L}_{+}$  and  $P_{1},\ldots ,P_{n}\in S_{m}^{+}$ , one has

$$
\det  \mathbb {E} _ {n} (L _ {1}, \dots , L _ {n}) = \exp \left(\frac {1}{n} \sum_ {i = 1} ^ {n} \log \det  L _ {i}\right) \tag {4.5}
$$
\end{corollary}

\paragraph{Background.}
-  $\operatorname{det}(X) = \prod_{j=1}^{m} X_{jj}$  for  $X \in \mathcal{L}$ .

These properties show that both  $\mathcal{L}$  and  $\mathcal{L}_{+}$  are closed under matrix addition and multiplication and that the operator  $\mathbb{D}$  interacts well with these operations.

Recall that  $S_{m}^{+}$  is defined as the collection of  $m \times m$  SPD matrices. We denote the space of  $m \times m$  symmetric space by  $S_{m}$ . Symmetric matrices and SPD matrices possess numerous algebraic and analytic properties that are well documented in [6]. Below are some of them to be used in what follows.

-  $\exp(S)$  is an SPD matrix for a symmetric matrix  $S$ , while  $\log(P)$  is a symmetric metric for an SPD matrix  $P$ .

- Diagonal elements of an SPD matrix are all positive. This can be seen from the fact that  $P_{jj} = e_j^\top Pe_j > 0$  for  $P \in S_m^+$ , where  $e_j$  is the unit vector with 1 at the  $j$ th coordinate and 0 elsewhere.

2.2. Cholesky decomposition. Cholesky decomposition, named after André-Louis Cholesky, represents a real  $m \times m$  SPD matrix  $P$  as a product of a lower triangular matrix  $L$  and its transpose, i.e.,  $P = LL^{\top}$ . If the diagonal elements of  $L$  are restricted to be positive, then the decomposition is unique according to Theorem 4.2.5 of [16]. Such lower triangular matrix, denoted by  $\mathcal{L}(P)$ , is called the Cholesky factor of  $P$ . Since in addition  $L = \mathcal{L}(LL^{\top})$  for each  $L \in \mathcal{L}_{+}$ , the map  $\mathcal{L}: S_{m}^{+} \to \mathcal{L}_{+}$  is bijective. In other words, there is one-to-one correspondence between SPD matrices and lower triangular matrices whose diagonal elements are all positive.

inverse of  $L$ , denoted by  $L_{\odot}^{-1}$ , is  $\mathbb{D}(L)^{-1} - \lfloor L\rfloor$ . The left translation by  $A\in \mathcal{L}_{+}$  is denoted by  $\ell_A:B\mapsto A\odot B$ . One can check that the differential of this operation at  $L\in \mathcal{L}_{+}$  is

$$
D _ {L} \ell_ {A}: X \mapsto \lfloor X \rfloor + \mathbb {D} (A) \mathbb {D} (X), \tag {3.2}
$$

where it is noted that the differential  $D_{L}\ell_{A}$  does not depend on  $L$ . Given the above expression, one can find that

$$
D _ {L} \ell_ {A}: X \mapsto \lfloor X \rfloor + \mathbb {D} (A) \mathbb {D} (X), \tag {3.2}
$$

where it is noted that the differential  $D_{L}\ell_{A}$  does not depend on  $L$ . Given the above expression, one can find that

$$
\begin{array}{l} \tilde {g} _ {A \odot L} ((D _ {L} \ell_ {A}) (X), (D _ {L} \ell_ {A}) (Y)) \\ = \tilde {g} _ {A \odot L} (\left\lfloor X \right\rfloor + \mathbb {D} (A) \mathbb {D} (X), \left\lfloor Y \right\rfloor + \mathbb {D} (A) \mathbb {D} (Y)) \\ = \tilde {g} _ {L} (X, Y). \\ \end{array}
$$


\begin{corollary}[Corollary 13]
Corollary 13 also suggests that the determinant of the Log-Cholesky average is equal to the determinant of the Log-Euclidean and affine-invariant averages. Thus,
\end{corollary}

\end{document}